% ============================================================
%  USPSC-Cap3-Metodologia.tex
%  Capítulo 3 — Metodologia
%  TCC: Análise de Sentimento de Publicações do X/Twitter
%       no Mercado Financeiro Brasileiro com FinBERT-PT-BR
% ============================================================
\chapter{Metodologia} \label{Cap:Metodologia}

Este capítulo descreve os procedimentos metodológicos adotados para o desenvolvimento da pesquisa. A construção do \textit{pipeline} de Processamento de Linguagem Natural (PLN) proposto demandou decisões sucessivas acerca do tipo de pesquisa, da fonte e do processo de coleta de dados, das etapas de tratamento e pré-processamento textual, das abordagens de classificação de sentimento e dos critérios de avaliação e interpretação dos resultados. Cada escolha é apresentada e justificada à luz dos objetivos
estabelecidos no \autoref{Cap:Introdução} e do referencial teórico discutido no \autoref{Cap:Fundamentacao_Teorica}.

% ============================================================
%  USPSC-Cap3-Metodologia_3.1-Caracterizacao_da_Pesquisa.tex
%  Seção 3.1 — Caracterização da Pesquisa
% ============================================================

\section{Caracterização da Pesquisa} \label{sec:caracterizacao_pesquisa}

A presente pesquisa pode ser caracterizada segundo quatro dimensões complementares: sua natureza, seus objetivos, a abordagem ao problema e os procedimentos técnicos empregados.

Quanto à sua \textbf{natureza}, trata-se de uma pesquisa \textbf{aplicada}, uma vez que visa produzir conhecimento direcionado à resolução de um problema concreto — a ausência de protocolos metodológicos replicáveis para análise de sentimento no domínio financeiro em língua portuguesa — com potencial de aplicação prática por analistas, gestores e pesquisadores da área \cite{galdi2018}.

Quanto aos \textbf{objetivos}, a pesquisa é \textbf{descritiva e exploratória}. Descritiva porque busca caracterizar a distribuição de polaridade e os padrões semânticos presentes no corpus de publicações coletadas, identificando regularidades temporais e temáticas sem estabelecer relações de causalidade. Exploratória porque investiga um domínio ainda
pouco estudado na literatura nacional — a análise de sentimento de publicações financeiras em português brasileiro via modelos de linguagem pré-treinados — e porque a avaliação comparativa das abordagens de classificação gera hipóteses para pesquisas futuras \cite{santos2023finbert}.

Quanto à \textbf{abordagem ao problema}, trata-se de uma pesquisa \textbf{quantitativa}. O corpus textual coletado é tratado computacionalmente para gerar representações numéricas de polaridade (positiva, negativa ou neutra) que, agregadas, produzem um índice de sentimento passível de análise estatística descritiva. Embora a interpretação dos resultados recorra ao contexto qualitativo de eventos macroeconômicos do período, a unidade central de análise é de natureza quantitativa \cite{kearney2014}.

Quanto aos \textbf{procedimentos técnicos}, a pesquisa adota o modelo \textbf{experimental-computacional}, caracterizado pela construção, aplicação e avaliação de um \textit{pipeline} de PLN sobre dados coletados programaticamente via interface de programação de aplicações (API). Esse modelo é consistente com a tradição metodológica da área de PLN aplicado a finanças, na qual a replicabilidade e a documentação pormenorizada dos procedimentos são requisitos centrais para a validade científica do trabalho
\citeonline{santos2023finbert}.

O \autoref{tab:resumo_pesquisa} sintetiza as escolhas metodológicas adotadas e suas respectivas justificativas.

\begin{table}[htb]
	\caption{Síntese da caracterização metodológica da pesquisa}
	\label{tab:resumo_pesquisa}
	\begin{tabular}{p{3.5cm}p{3.5cm}p{8cm}}
		\hline
		\textbf{Dimensão} & \textbf{Classificação} & \textbf{Justificativa} \\
		\hline
		Natureza & Aplicada & Geração de protocolo metodológico replicável para o domínio financeiro em PT-BR \\
		\hline
		Objetivos & Descritiva e exploratória & Caracterização de padrões de sentimento; investigação de domínio pouco estudado na literatura nacional \\
		\hline
		Abordagem ao problema & Quantitativa & Mensuração numérica de polaridade via PLN e construção de índice de sentimento \\
		\hline
		Procedimentos técnicos & Experimental-computacional & Construção, aplicação e avaliação de \textit{pipeline} de PLN sobre dados coletados via API \\
		\hline
	\end{tabular}
	\legend{Fonte: Elaborado pelo autor.}
\end{table}

% ============================================================
%  USPSC-Cap3-Metodologia_3.2-Fonte_de_Dados_e_Coleta.tex
%  Seção 3.2 — Fonte de Dados e Processo de Coleta
% ============================================================

\section{Fonte de Dados e Processo de Coleta}
\label{sec:fonte_dados}

\subsection{Justificativa da Escolha do X/Twitter como Fonte}
\label{subsec:justificativa_twitter}

O X/Twitter foi selecionado como fonte primária de dados por congregar características estruturais que o tornam especialmente adequado para a análise de sentimento aplicada ao mercado financeiro. Em primeiro lugar, a plataforma opera como canal de comunicação em tempo real, com publicações cronológicas e públicas que registram o humor e as expectativas dos agentes financeiros de forma praticamente instantânea \cite{bollen2011}. Em segundo lugar, a arquitetura da plataforma — baseada em publicações curtas (\textit{tweets}) com limite de caracteres — produz textos concisos e com
orientação semântica, mais adequados à tarefa de classificação de polaridade do que textos longos de natureza mais ambígua \cite{ranco2015}.

A literatura internacional consolidou o X/Twitter como objeto de pesquisa em finanças, desde os trabalhos pioneiros de \citeonline{bollen2011} e \citeonline{antweiler2004}, que demonstraram a relação entre o conteúdo das publicações e o comportamento de mercado. No contexto brasileiro, trabalhos como os de \citeonline{sousa2025} e \citeonline{falcao2020} também recorreram à plataforma como fonte de dados para análise de sentimento financeiro.

Em comparação com alternativas como artigos completos de veículos de imprensa, as publicações do X/Twitter apresentam duas vantagens metodológicas relevantes: (i) maior facilidade de extração estruturada de dados via API, dispensando técnicas complexas de \textit{web scraping}; e (ii) formato textual conciso e objetivo, mais adequado para análise de polaridade semântica no nível de documento.

\subsection{Veículos Selecionados}
\label{subsec:veiculos}

Como \textit{proxy} para o sentimento do mercado financeiro brasileiro, foram selecionadas as publicações de dois veículos especializados em economia e finanças com presença ativa na plataforma X/Twitter: \textbf{InfoMoney} (\texttt{@InfoMoney}) e \textbf{Investing Brasil} (\texttt{@InvestingBrasil}).

A escolha desses veículos é fundamentada em três critérios:

\begin{enumerate}
	\item \textbf{Relevância editorial}: ambos são referências reconhecidas de jornalismo econômico-financeiro no Brasil, com cobertura sistemática de eventos macroeconômicos, decisões de política monetária, resultados corporativos e movimentações do mercado de capitais nacional;

	\item \textbf{Volume e regularidade de publicações}: as contas institucionais mantêm frequência elevada e contínua de publicações, garantindo densidade suficiente de dados para a construção de séries temporais de sentimento;

	\item \textbf{Consistência do público-alvo}: ao contrário de contas de pessoas físicas, cujo conteúdo pode ser heterogêneo, os veículos especializados produzem publicações com foco temático predominantemente financeiro, reduzindo a necessidade de filtragem por relevância.
\end{enumerate}

Reconhece-se que a opção por veículos institucionais em detrimento de contas de investidores individuais representa uma escolha metodológica deliberada, que delimita o escopo do sentimento capturado ao discurso editorial — e não ao sentimento difuso dos participantes de varejo do mercado. Essa limitação circunscreve o alcance das generalizações admissíveis e é reconhecida como restrição metodológica deste trabalho.

\subsection{Período de Análise}
\label{subsec:periodo}

O corpus compreende publicações realizadas entre \textbf{17 de outubro de 2025} e \textbf{12 de fevereiro de 2026}, totalizando aproximadamente quatro meses de dados — intervalo determinado pelo registro mais antigo e pelo mais recente efetivamente presentes na base coletada via API. Esse recorte temporal foi determinado por um conjunto de restrições metodológicas e operacionais.

Do ponto de vista \textbf{técnico}, o acesso à API do X/Twitter está condicionado a planos comerciais com janelas históricas limitadas. A recuperação retroativa de grandes volumes de dados requer subscrição ao plano \textit{Pro} ou \textit{Enterprise} da API, cujos custos são proibitivos para pesquisas acadêmicas sem financiamento dedicado. O período analisado corresponde à janela de dados coletados em tempo real durante a execução da pesquisa, garantindo acesso completo e sem truncamento ao histórico de publicações.

Do ponto de vista \textbf{analítico}, o período de aproximadamente quatro meses abrange um conjunto relevante de eventos macroeconômicos do cenário brasileiro — incluindo reuniões do Comitê de Política Monetária (COPOM), divulgação de indicadores de inflação (IPCA e INPC), dados de atividade econômica e movimentações relevantes nos preços dos ativos — o que possibilita a avaliação qualitativa da coerência entre os índices de sentimento extraídos e o contexto econômico-financeiro do período.

\subsection{Processo de Coleta via API}
\label{subsec:coleta_api}

\subsubsection{\textit{Endpoint} Utilizado e Autenticação}
\label{subsubsec:endpoint}

A coleta foi realizada por meio do \textit{endpoint} \textbf{GET \texttt{/2/users/\{id\}/tweets}} da API v2 do X/Twitter\footnote{Documentação oficial: \url{https://developer.twitter.com/en/docs/twitter-api/tweets/timelines/api-reference/get-users-id-tweets}. Acesso em: 12 fev. 2026.}, que retorna a linha do tempo de publicações de um usuário específico identificado pelo seu \texttt{id} numérico. A opção por esse \textit{endpoint} — em vez do \textit{endpoint} de busca (\textit{Search}) — é metodologicamente justificada pelo escopo da coleta: como os veículos de interesse (InfoMoney e Investing Brasil) são previamente conhecidos e possuem identificadores únicos fixos, a recuperação direta pela linha do tempo do usuário é mais eficiente e determinística do que uma busca por palavras-chave, eliminando a necessidade de filtros adicionais de autoria.

A autenticação foi realizada via \textbf{Bearer Token} (\textit{app-only authentication}), mecanismo que concede acesso de leitura ao conteúdo público da plataforma sem requerer autorização individual de cada conta coletada. O token foi armazenado como variável de ambiente (\texttt{TWITTER\_ACCESS\_TOKEN}) e transmitido no cabeçalho HTTP de cada requisição no formato \texttt{Authorization: Bearer <token>}.

\subsubsection{Parâmetros da Requisição}
\label{subsubsec:parametros_api}

Para cada veículo monitorado, as requisições foram parametrizadas conforme o \autoref{qua:params_api}. O intervalo temporal foi delimitado pelos campos \texttt{start\_time} e \texttt{end\_time}, ambos no formato ISO~8601 (\texttt{YYYY-MM-DDTHH:mm:ssZ}, UTC). O volume máximo de resultados por página foi fixado em 100 registros — limite superior permitido pelo \textit{endpoint} —, e a paginação foi controlada pelo parâmetro \texttt{pagination\_token}, utilizado de forma iterativa até o esgotamento dos resultados disponíveis.

\begin{quadro}[htb]
	\caption{\label{qua:params_api}Parâmetros utilizados nas requisições ao \textit{endpoint} \texttt{GET /2/users/\{id\}/tweets}}
	
	\begin{tabular}{|p{4.5cm}|p{10.5cm}|}
		\hline
		\textbf{Parâmetro} & \textbf{Valor / Descrição} \\
		\hline
		
		\texttt{id} (path) & Identificador numérico do veículo: \texttt{@InfoMoney} e \texttt{@InvestingBrasil}, passados dinamicamente a cada ciclo de coleta \\
		\hline
		
		\texttt{start\_time} & Limite inferior do intervalo temporal (ISO 8601, UTC) \\
		\hline
		
		\texttt{end\_time} & Limite superior do intervalo temporal (ISO 8601, UTC) \\
		\hline
		
		\texttt{max\_results} & \texttt{100} (máximo permitido pelo \textit{endpoint}) \\
		\hline
		
		\texttt{tweet.fields} & \texttt{created\_at}, \texttt{note\_tweet}, \texttt{author\_id}, \texttt{public\_metrics}, \texttt{lang}, \texttt{source}, \texttt{entities}, \texttt{context\_annotations}, \texttt{geo} \\
		\hline
		
		\texttt{user.fields} & \texttt{name}, \texttt{username}, \texttt{location}, \texttt{description}, \texttt{public\_metrics} \\
		\hline
		
		\texttt{pagination\_token} & Token retornado pela API para recuperação da próxima página de resultados; omitido na primeira requisição de cada ciclo \\
		\hline
	\end{tabular}
	\legend{Fonte: Elaborado pelo autor com base na documentação da API v2 do X/Twitter e no código-fonte do \textit{pipeline} de coleta.}
\end{quadro}

Um aspecto técnico relevante diz respeito ao campo \texttt{note\_tweet}: publicações de usuários com assinatura X Premium podem ultrapassar o limite convencional de 280 caracteres por meio desse campo, que armazena o texto estendido separadamente do campo \texttt{text} padrão. O \textit{pipeline} implementado prioriza o conteúdo de \texttt{note\_tweet} quando disponível, recorrendo ao \texttt{text} convencional apenas na sua ausência — decisão que garante a recuperação do conteúdo textual integral das publicações. Registros sem conteúdo textual em nenhum dos dois campos são descartados durante a coleta.

\subsubsection{\textit{Pipeline} de Coleta e Armazenamento}
\label{subsubsec:pipeline_coleta}

O \textit{pipeline} de coleta é a primeira etapa (\texttt{pipeline/01\_extraction/}) de uma estrutura modular de seis estágios numerados. Os módulos que compõem essa etapa, com suas
responsabilidades bem delimitadas, são:

\begin{itemize}
	\item \texttt{app/core/extraction/twitter\_api.py}: implementa a classe \texttt{TwitterAPIExtractor}, que herda da classe abstrata \texttt{BaseExtractor} e é responsável pela construção e execução
	das requisições HTTP à API via biblioteca \texttt{requests}. Recebe como parâmetros o identificador do usuário (\texttt{x\_user\_id}), o intervalo temporal e, opcionalmente, o token de paginação;
	
	\item \texttt{app/core/extraction/base\_extractor.py}: define a classe abstrata \texttt{BaseExtractor} (ABC), que estabelece o contrato de interface para qualquer extrator de dados do 	\textit{pipeline}, garantindo extensibilidade para futuras fontes de dados além do X/Twitter;
	
	\item \texttt{app/shared/schemas.py}: módulo transversal a todas as etapas do \textit{pipeline}, contendo as \textit{dataclasses} de \textit{parsing} da resposta JSON da API — \texttt{Tweet}, \texttt{Tweets}, \texttt{NoteTweet}, \texttt{PublicMetrics}, \texttt{Entities}, \texttt{Hashtag} e \texttt{Meta} —, que estruturam os campos retornados em objetos Python tipados;
	
	\item \texttt{app/core/extraction/collection\_log.py}: define a classe \texttt{SearchTerm}, que representa um registro de coleta agendada com os campos \texttt{x\_user\_id}, 	\texttt{from\_date\_time} e \texttt{to\_date\_time}, e gerencia a inserção dos termos de busca na tabela \texttt{collection\_log};
	
	\item \texttt{app/shared/db/database.py}: módulo transversal que implementa a classe \texttt{DatabaseManager}, responsável pela interface de comunicação com o banco de dados PostgreSQL — incluindo as operações de consulta, inserção e atualização utilizadas por todos os estágios do \textit{pipeline};
	
	\item \texttt{app/core/extraction/\_\_main\_\_.py}: ponto de entrada do estágio de coleta (invocado via \texttt{python -m app.core.extraction}, que orquestra o ciclo completo: recupera
	registros com status \texttt{pending} em \texttt{collection\_log}, executa a paginação iterativa para cada um deles e persiste os resultados no banco de dados.
\end{itemize}

Os dados coletados foram persistidos em um banco de dados \textbf{PostgreSQL~15}, provisionado em contêiner \textbf{Docker} para garantir portabilidade e isolamento do ambiente. A infraestrutura é inicializada via \texttt{docker-compose.yml}, que sobe automaticamente o PostgreSQL (porta~5432) e o pgAdmin~4 (porta~5050), executando os scripts \texttt{infra/01-init-database.sh} — criação de usuário e permissões — e \texttt{infra/02-create-tables.sql} — criação das tabelas e índices. O esquema do banco é composto por três tabelas, cujas estruturas são descritas no \autoref{qua:schema_db}.

\begin{quadro}[!htb]
	\caption{\label{qua:schema_db}Estrutura das tabelas do banco de dados PostgreSQL}
	
	\begin{tabular}{|p{4.5cm}|p{4.5cm}|p{5.8cm}|}
		\hline
		\textbf{Tabela} & \textbf{Campo} & \textbf{Descrição} \\
		\hline
		\multirow{8}{*}{\texttt{tweets}} & \texttt{tweet\_id} & Identificador único da publicação na plataforma (\texttt{UNIQUE}) \\ & \texttt{username} & Identificador numérico do veículo autor \\
		& \texttt{note\_tweet} & Conteúdo textual utilizado na análise (\texttt{note\_tweet.text} ou \texttt{text}) \\ & \texttt{created\_at} & Data e hora de publicação (com fuso horário) \\ & \texttt{likes}
		& Número de curtidas \\ & \texttt{hashtags} & Lista de \textit{hashtags} (JSONB) \\ & \texttt{tweet} & Objeto JSON completo retornado pela API
		(JSONB) \\
		
		\hline
		\multirow{7}{*}{\texttt{tweets\_classification}} & \texttt{tweet\_id} & FK → \texttt{tweets.id} (\texttt{ON DELETE CASCADE}) \\ & \texttt{sentiment} & Polaridade atribuída: \textit{positivo}, \textit{negativo} ou \textit{neutro} \\ & \texttt{why\_sentiment} & Justificativa textual para a polaridade classificada \\ & \texttt{is\_finance\_news}   & Relevância financeira classificada: 1~=~financeiro, 0~=~não financeiro \\ & \texttt{why\_is\_finance\_news} & Justificativa textual para a classificação de relevância \\ & \texttt{classificator} & Identificador do classificador: \texttt{"Humano"}, \texttt{"FinBERT-PT-BR"}, \texttt{"BERTimbau"} ou outro modelo futuro \\
		
		\hline
		\multirow{6}{*}{\texttt{collection\_log}} & \texttt{search\_term} & Parâmetros da coleta agendada (JSONB) \\ & \texttt{tweets\_collected}    & Total de registros inseridos no 	ciclo \\
		& \texttt{start\_time} & Início da execução do ciclo de coleta \\ & \texttt{end\_time} & Término da execução do ciclo de coleta \\ 	& \texttt{status} & Estado do ciclo: \texttt{pending}, 		\texttt{completed} ou \texttt{partially\_completed} \\ & \texttt{error\_message} & Mensagem de erro, se houver \\
		\hline
	\end{tabular}
	\legend{Fonte: Elaborado pelo autor com base no schema \texttt{infra/02-create-tables.sql} do repositório do \textit{pipeline}.}
\end{quadro}

A tabela \texttt{collection\_log} cumpre função de controle de idempotência: cada combinação de veículo e intervalo temporal é registrada com status \texttt{pending} antes da execução, sendo
atualizada para \texttt{completed} ao final de um ciclo sem falhas de inserção, ou para \texttt{partially\_completed} caso algum registro não tenha sido persistido. Esse mecanismo permite a retomada segura de coletas interrompidas e a auditoria do histórico de execuções sem risco de duplicação de dados, uma vez que a coluna \texttt{tweet\_id} na tabela \texttt{tweets} é definida com restrição \texttt{UNIQUE}.

% ============================================================
%  USPSC-Cap3-Metodologia_3.3-Rotulagem_Manual.tex
%  Seção 3.3 — Rotulagem Manual do Corpus
% ============================================================

\section{Rotulagem Manual do Corpus}
\label{sec:rotulagem}

A rotulagem manual constitui etapa fundamental para a construção do conjunto de referência (\textit{gold standard}) utilizado na avaliação das abordagens de classificação de sentimento descritas na \autoref{sec:abordagens}. O processo foi conduzido pelo próprio pesquisador por meio de uma ferramenta de anotação desenvolvida especificamente para este trabalho, implementada em Python com a biblioteca \texttt{streamlit} e integrada ao banco de dados PostgreSQL descrito na \autoref{subsubsec:pipeline_coleta}.

\subsection{Protocolo de Anotação em Duas Etapas}
\label{subsec:protocolo_anotacao}

O protocolo de anotação foi estruturado em duas etapas sequenciais, aplicadas a cada publicação individualmente:

\textbf{Etapa 1 — Relevância financeira.} O anotador classifica a publicação em uma de duas categorias mutuamente exclusivas: \textit{tweet financeiro} (valor~1), quando o conteúdo se relaciona diretamente ao mercado financeiro, à economia, a investimentos ou a eventos com impacto sobre ativos financeiros; ou \textit{não financeiro} (valor~0), quando o conteúdo é de outra natureza — editorial, esportivo, de entretenimento ou sem relação com o contexto econômico-financeiro. Apenas as publicações classificadas como financeiras prosseguem para a segunda etapa.

\textbf{Etapa 2 — Polaridade de sentimento.} Para as publicações identificadas como financeiras, o anotador atribui uma das três classes de polaridade:

\begin{itemize}
	\item \textbf{Positivo}: o conteúdo veicula uma perspectiva favorável ao mercado ou a agentes econômicos — crescimento, resultados acima do esperado, alívio de incertezas, aprovação de
	medidas de estímulo, entre outros;
	
	\item \textbf{Neutro}: o conteúdo tem caráter predominantemente informativo e factual, sem atribuir valoração explícita ao evento relatado;

	\item \textbf{Negativo}: o conteúdo veicula uma perspectiva desfavorável — queda de indicadores, resultados abaixo do esperado, aumento de incerteza, crises ou eventos de risco sistêmico.
\end{itemize}

A separação entre relevância financeira e polaridade de sentimento em etapas distintas é uma decisão metodológica deliberada: ela evita que publicações de conteúdo não financeiro --- mesmo que textualmente negativas ou positivas --- contaminem a análise de sentimento do mercado financeiro.

\subsection{Ferramenta de Anotação}
\label{subsec:ferramenta_anotacao}

A ferramenta de anotação constitui o segundo estágio do \textit{pipeline} (\texttt{app/dashboard/}) e é composta por dois arquivos:

\begin{itemize}
	\item \texttt{app/dashboard/app.py}: \textit{launcher} multi-página Streamlit, invocado via \texttt{streamlit run app/dashboard/app.py} (acessível em \texttt{http://localhost:8501}). Agrega em uma única interface 6 páginas:
	
	\begin{itemize}
		\item \textit{Anotação} (\texttt{pages/annotation.py});
		\item \textit{Eploração de Dados} (\texttt{pages/exploration.py});
		\item \textit{Split do Dadaset} (\texttt{pages/dataset\_split.py});
		\item \textit{Fine-tuning BERTimbau} (\texttt{pages/fine\_tuning.py});
		\item \textit{Classificação} (\texttt{pages/classification.py}) e;
		\item \textit{Avaliação} (\texttt{pages/evaluation.py}.
	\end{itemize}

\end{itemize}

A interface de classificação é composta por três componentes principais:

\begin{enumerate}
	\item \textbf{Exibição do conteúdo}: para cada publicação, a ferramenta exibe o texto integral do campo \texttt{note\_tweet} e a data de publicação, permitindo ao anotador contextualizar
	temporalmente o conteúdo antes de classificá-lo;
	
	\item \textbf{Botões de classificação}: dois botões para a decisão de relevância financeira (Etapa~1) e, condicionalmente, três botões para a atribuição de polaridade (Etapa~2). O estado da classificação corrente é destacado visualmente, permitindo revisão e correção a qualquer momento;
	
	\item \textbf{Registro de justificativa}: um formulário opcional (\textit{dialog} Streamlit) permite que o anotador registre, em texto livre, a justificativa para cada decisão — tanto para a relevância financeira quanto para a polaridade atribuída. Esse registro é armazenado nos campos \texttt{why\_is\_finance\_news} e \texttt{why\_sentiment} da tabela \texttt{tweets\_classification}, enriquecendo o \textit{gold standard} com raciocínio explícito passível de auditoria futura.
\end{enumerate}

A navegação pelo corpus é sequencial, com barra de progresso indicando o percentual de publicações já classificadas. A ferramenta prioriza automaticamente as publicações ainda sem classificação
humana, identificadas pelo campo \texttt{has\_human\_classification}. Ao final de cada publicação, um histórico tabular exibe todas as classificações anteriores registradas em \texttt{tweets\_classification}, com os campos \textit{classificador}, \textit{sentimento}, \textit{relevância financeira} e as respectivas justificativas.

\subsection{Armazenamento das Anotações}
\label{subsec:armazenamento_anotacoes}

Cada decisão de classificação é persistida em duas estruturas complementares do banco de dados:

\begin{itemize}
	\item \textbf{Tabela \texttt{tweets}}: os campos \texttt{is\_finance\_news} e \texttt{sentiment} são atualizados diretamente, tornando a classificação imediatamente disponível para as etapas de análise subsequentes;
	
	\item \textbf{Tabela \texttt{tweets\_classification}}: registra o histórico completo de cada decisão com o identificador do classificador (\texttt{classificator = ``Humano''}), as justificativas textuais e os valores atribuídos, garantindo rastreabilidade e possibilitando análises de concordância em estudos futuros.
\end{itemize}


\subsection{Particionamento do Corpus Rotulado para Treinamento e Avaliação}
\label{subsec:particionamento}

O corpus rotulado manualmente cumpre, neste trabalho, dois propósitos metodológicos distintos: (i) constitui o conjunto de referência (\textit{gold standard}) utilizado na avaliação comparativa das abordagens de classificação descritas na \autoref{sec:abordagens}; e (ii) é a fonte de exemplos supervisionados para o \textit{fine-tuning} de modelos pré-treinados ajustados localmente, como o BERTimbau (\autoref{subsec:bertimbau}). Esse duplo uso impõe uma decisão de protocolo: para que a comparação entre abordagens que viram os rótulos durante o treinamento (e.g., BERTimbau ajustado) e abordagens que não os viram (e.g., FinBERT-PT-BR e métodos léxicos) seja metodologicamente honesta, todas as abordagens devem ser avaliadas sobre o \textbf{mesmo subconjunto de teste}, definido a priori, antes de qualquer treinamento \cite{jurafsky2026,caseli2024cap1}.

O particionamento é gerenciado pela classe \texttt{DatasetSplitRepository} (\texttt{app/shared /db/dataset\_split.py}), que implementa o método \texttt{assign\_split()} e persiste o resultado na tabela \texttt{dataset\_split}. A operação é idempotente: uma vez realizada, é reutilizada por todas as execuções subsequentes do \textit{pipeline}; o parâmetro \texttt{force=True} permite regenerá-la caso o corpus rotulado seja ampliado. A \autoref{tab:dataset_split} descreve o esquema da tabela.

\begin{table}[!htbp]
	\caption{Esquema da tabela \texttt{dataset\_split}}
	\label{tab:dataset_split}
	\centering
	\begin{tabular}{llp{8.2cm}}
		\hline
		\textbf{Coluna} & \textbf{Tipo} & \textbf{Descrição} \\
		\hline
		\texttt{tweet\_id} & INTEGER & Chave estrangeira para \texttt{tweets.id}; único na tabela \\
		\hline
		\texttt{split}     & VARCHAR & \texttt{'train'} ou \texttt{'test'} \\
		\hline
		\texttt{fold}      & SMALLINT & Fold do K-Fold ($1$--$4$) quando \texttt{split='train'}; \texttt{NULL} quando \texttt{split='test'} \\
		\hline
	\end{tabular}
	\legend{Fonte: Elaborado pelo autor.}
\end{table}

O procedimento compreende duas etapas sequenciais, ambas com estratificação pela classe de polaridade e semente aleatória fixa (\texttt{random\_state=42}), que preserva a distribuição original das três classes em cada conjunto — condição particularmente relevante dado o desbalanceamento esperado em favor da classe neutra, característico de publicações jornalísticas de natureza factual \cite{loughran2011}:

\begin{enumerate}
	\item \textbf{Hold-out de teste} ($\approx$~15\%): por meio de \texttt{train\_test\_split} com estratificação pelo sentimento, extrai-se um subconjunto fixo reservado exclusivamente para a avaliação comparativa final. Os registros recebem \texttt{split=`test'} e \texttt{fold=NULL}. Nenhum classificador — incluindo o BERTimbau — tem acesso a esses registros durante o treinamento, garantindo que as estimativas de desempenho reportadas no \autoref{Cap:Ava_Exp} sejam livres de viés de seleção;
	
	\item \textbf{Conjunto de treinamento} ($\approx$~85\%): os registros restantes recebem \texttt{split=`train'} e são subdivididos em $K=4$ folds por validação cruzada estratificada (\texttt{StratifiedK Fold}, \texttt{shuffle=True}, \texttt{random\_state=42}), cada um marcado com \texttt{fold}~$\in \{1, 2, 3, 4\}$. Esse conjunto é utilizado exclusivamente pelo protocolo de \textit{fine-tuning} do BERTimbau, detalhado na \autoref{subsubsec:arquitetura_bertimbau}; os demais classificadores (OpLexicon, SentiLex-PT e FinBERT-PT-BR) não requerem dados de treinamento e, portanto, não dependem dessa subdivisão.
\end{enumerate}

A \autoref{fig:dataset_split} ilustra a estrutura resultante do particionamento.

\begin{figure}[!htb]
	\centering
	\caption{\label{fig:dataset_split}Estrutura do particionamento estratificado do corpus rotulado}
	\resizebox{0.82\textwidth}{!}{%
		\begin{tikzpicture}
			[
				box/.style   = {rectangle, rounded corners=3pt, draw=black,fill=gray!12, align=center, minimum height=1.0cm,font=\small},
				fold/.style  = {rectangle, rounded corners=3pt, draw=black!60, fill=gray!5,  align=center, minimum height=0.9cm, font=\small},
				lbl/.style   = {font=\scriptsize, align=center},
				node distance = 0.35cm and 0.4cm
			]
			
			%% Corpus total
			\node[box, text width=11.5cm, minimum height=1.1cm]
			(corpus){Corpus rotulado (tweets financeiros com anotação humana)};
			
			%% Hold-out
			\node[box, text width=3.0cm, below left=0.9cm and 1.8cm of corpus,
			fill=gray!25] (test)
			{Hold-out\\$\approx$~15\%\\\texttt{split=`test'}};
			
			%% Train pool
			\node[box, text width=7.8cm, below right=0.9cm and -1.0cm of corpus]
			(train) {Conjunto de treinamento~$\approx$~85\%\quad\texttt{split=`train'}};
			
			%% Folds
			\node[fold, text width=1.55cm, below=0.6cm of train, xshift=-2.7cm]
			(f1) {Fold 1\\\texttt{fold=1}};
			\node[fold, text width=1.55cm, right=0.35cm of f1] (f2)
			{Fold 2\\\texttt{fold=2}};
			\node[fold, text width=1.55cm, right=0.35cm of f2] (f3)
			{Fold 3\\\texttt{fold=3}};
			\node[fold, text width=1.55cm, right=0.35cm of f3] (f4)
			{Fold 4\\\texttt{fold=4}};
			
			%% Arrows
			\draw[->, thick] (corpus.south) -- ++(0,-0.4) -| (test.north);
			\draw[->, thick] (corpus.south) -- ++(0,-0.4) -| (train.north);
			\draw[->, thick] (train.south)  -- ++(0,-0.3) -| (f1.north);
			\draw[->, thick] (train.south)  -- ++(0,-0.3) -| (f2.north);
			\draw[->, thick] (train.south)  -- ++(0,-0.3) -| (f3.north);
			\draw[->, thick] (train.south)  -- ++(0,-0.3) -| (f4.north);
			
			%% Labels
			\node[lbl, above right=0.05cm and 0.1cm of test.north east]
			{\texttt{random\_state=42}\\estratificado por sentimento};
			\node[lbl, below=0.15cm of f2.south, xshift=0.95cm]
			{\texttt{StratifiedKFold}$(K=4,\;\texttt{shuffle=True},\;\texttt{random\_state=42})$};
			
		\end{tikzpicture}%
	}
	\legend{Fonte: Elaborado pelo autor.}
\end{figure}

A fixação da semente aleatória e a persistência dos identificadores na tabela \texttt{dataset\_split} asseguram a reprodutibilidade integral do protocolo: qualquer execução subsequente do \textit{pipeline} sobre o mesmo corpus rotulado recupera os mesmos conjuntos, condição necessária para que comparações entre abordagens introduzidas em momentos distintos da pesquisa permaneçam metodologicamente consistentes.


\subsection{Limitações do Processo de Rotulagem}
\label{subsec:lim_rotulagem}

O processo de rotulagem foi conduzido por um único anotador — o próprio pesquisador —, o que inviabiliza o cálculo de métricas de concordância interanotadores (e.g., Kappa de Cohen). A subjetividade inerente à tarefa, em especial para publicações de polaridade ambígua ou contextualmente dependente, representa uma fonte de incerteza não quantificada nos resultados, reconhecida como restrição metodológica deste trabalho.

% ============================================================
%  USPSC-Cap3-Metodologia_3.4-Analise_Exploratoria_dos_Dados.tex
%  Seção 3.4 — Análise Exploratória dos Dados
% ============================================================

\section{Análise Exploratória dos Dados}
\label{sec:eda}

Após a coleta e persistência dos registros no banco de dados PostgreSQL, o corpus foi submetido a uma Análise Exploratória dos Dados (AED) com três objetivos: (i)~caracterizar a dimensão e a composição do conjunto coletado; (ii)~verificar a completude dos campos e identificar valores ausentes; e (iii)~examinar a distribuição do comprimento dos textos e detectar publicações atípicas que pudessem comprometer a etapa de classificação de sentimento. A AED foi implementada em Python como o terceiro estágio do \textit{pipeline} (\texttt{app/dashboard/}), estruturado em dois módulos Streamlit com responsabilidades complementares, ambos acessados via \texttt{streamlit run app/dashboard/app.py}:

\begin{itemize}
	\item \texttt{app/dashboard/pages/exploration.py}: responsável pela análise textual do corpus — dimensões do conjunto, completude dos campos, distribuição do comprimento dos textos (caracteres e tokens) e
	identificação de publicações atípicas;
	
	\item \texttt{app/dashboard/pages/eda.py}: responsável pela visualização da distribuição de sentimentos e de relevância financeira entre as publicações já rotuladas, por meio de gráficos de setores interativos.
\end{itemize}

Ambos os módulos obtêm os dados diretamente do banco de dados PostgreSQL via \texttt{app/shared/db/database.py}, com o auxílio das bibliotecas \texttt{pandas}, \texttt{plotly} e \texttt{streamlit}. Os resultados das subseções a seguir são provenientes do módulo \texttt{eda.py}.


\subsection{Dimensões do Corpus e Distribuição por Veículo}
\label{subsec:dimensoes}

O corpus final, exportado após o encerramento da coleta, é composto por \textbf{3.563 publicações}, distribuídas entre dois veículos monitorados conforme o \autoref{tab:dist_veiculo}. O InfoMoney responde pela ampla maioria das publicações (88,5\% do total), enquanto o InvestingBrasil contribui com os 11,5\% restantes. Essa assimetria reflete diferenças estruturais no ritmo de publicação das contas institucionais e deve ser considerada na etapa de construção do índice de sentimento agregado.

\begin{table}[!htbp]
	\centering
	\caption{Distribuição do corpus por veículo monitorado}
	\label{tab:dist_veiculo}
	\begin{tabular}{lrr}
		\hline
		\textbf{Veículo} & \textbf{Publicações} & \textbf{Participação (\%)} \\
		\hline
		InfoMoney        & 3.154 & 88,5 \\
		\hline
		InvestingBrasil  &   409 & 11,5 \\
		\hline
		\textbf{Total}   & \textbf{3.563} & \textbf{100,0} \\
	\end{tabular}
	\legend{Fonte: Elaborado pelo autor com base nos dados coletados.}
\end{table}

\begin{figure}[!htb]
	\caption{\label{fig:dist_corpus}Distribuição do corpus por veículo (esquerda) e distribuição temporal mensal (direita)}
	\begin{minipage}[t]{0.48\textwidth}
		\centering
		\includegraphics[width=\textwidth]{USPSC-TA-Textual/USPSC-Cap3-Metodologia/figuras/Dimensao_Conjunto.png}
		\legend{Fonte: Elaborado pelo autor.}
	\end{minipage}
	\hfill
	\begin{minipage}[t]{0.48\textwidth}
		\centering
		\includegraphics[width=\textwidth]{USPSC-TA-Textual/USPSC-Cap3-Metodologia/figuras/Ditribuicao_Temporal.png}
		\legend{Fonte: Elaborado pelo autor.}
	\end{minipage}
\end{figure}

A distribuição temporal das publicações ao longo do período analisado é apresentada no \autoref{tab:dist_mensal} e ilustrada na \autoref{fig:dist_corpus}. Observa-se que o InfoMoney não apresenta registros para outubro de 2025 — primeiro mês da janela de coleta —, ao passo que o InvestingBrasil possui 36 publicações nesse período. A ausência do InfoMoney em outubro deve ser investigada para determinar se decorreu de limitação da janela de coleta da API, de interrupção editorial do veículo ou de outra causa operacional. Os meses de dezembro de 2025 e janeiro de 2026 concentram o maior volume de publicações para ambos os veículos, e fevereiro de 2026 apresenta volume reduzido por corresponder a um período parcialmente coletado (até o dia~12).

\begin{table}[!htbp]
	\centering
	\caption{Distribuição mensal das publicações por veículo}
	\label{tab:dist_mensal}
	\begin{tabular}{lrrr}
		\hline
		\textbf{Mês}     & \textbf{InfoMoney} & \textbf{InvestingBrasil} & \textbf{Total} \\
		\hline
		Out/2025  &    0 &  36 &    36 \\
		\hline
		Nov/2025  &  723 &  94 &   817 \\
		\hline
		Dez/2025  &  997 & 119 & 1.116 \\
		\hline
		Jan/2026  & 1044 &  88 & 1.132 \\
		\hline
		Fev/2026  &  390 &  72 &   462 \\
		\hline
		\textbf{Total} & \textbf{3.154} & \textbf{409} & \textbf{3.563} \\
	\end{tabular}
	\legend{Fonte: Elaborado pelo autor com base nos dados coletados.}
\end{table}

 
\subsection{Completude dos Campos}
\label{subsec:completude}

A verificação de valores ausentes foi realizada sobre os campos diretamente coletados via API do X/Twitter. O \autoref{tab:missing} sumariza a completude dos campos analisados, e a \autoref{fig:completude} apresenta a distribuição percentual de valores ausentes de forma visual.

\begin{table}[!htbp]
	\centering
	\caption{Completude dos campos coletados via API}
	\label{tab:missing}
	\begin{tabular}{lrrl}
		\hline
		\textbf{Campo} & \textbf{Ausentes} & \textbf{\% Ausente} &
		\textbf{Observação} \\
		\hline
		\texttt{tweet\_id}   &   0 & 0,0 & Completo \\
		\hline
		\texttt{username}    &   0 & 0,0 & Completo \\
		\hline
		\texttt{note\_tweet} &   0 & 0,0 & Completo \\
		\hline
		\texttt{created\_at} &   0 & 0,0 & Completo \\
		\hline
		\texttt{likes}       &   0 & 0,0 & Completo \\
		\hline
		\texttt{hashtags}    & 163 & 4,6 & Publicações sem \textit{hashtag} \\
		\hline
	\end{tabular}
	\legend{Fonte: Elaborado pelo autor com base nos dados coletados.}
\end{table}

\begin{figure}[!htb]
	\centering
	\caption{\label{fig:completude}Percentual de valores ausentes nos campos
		coletados via API}
	\includegraphics[width=0.85\textwidth]{USPSC-TA-Textual/USPSC-Cap3-Metodologia/figuras/Completude_Campos.png}
	\legend{Fonte: Elaborado pelo autor.}
\end{figure}

O corpus apresenta completude integral em cinco dos seis campos analisados. A única exceção é o campo \texttt{hashtags} (163 registros; 4,6\%), o que não constitui problema metodológico: o uso de \textit{hashtags} é facultativo nas publicações e sua ausência não implica a inexistência de conteúdo textual classificável — o campo \texttt{note\_tweet}, unidade central de análise, está integralmente preenchido em todos os 3.563 registros do corpus.


\subsection{Distribuição do Comprimento dos Textos}
\label{subsec:comprimento}

Uma característica estrutural relevante do corpus, evidenciada pela AED, é o comprimento inusualmente elevado dos textos em relação ao formato convencional de publicações do X/Twitter. A mediana de \textbf{778	caracteres} e de \textbf{121 tokens} por publicação — valores muito superiores ao limite padrão de 280 caracteres da plataforma — confirma que a quase totalidade das publicações foi veiculada por meio do campo \texttt{note\_tweet} (texto estendido), disponível para usuários com assinatura X Premium. O \autoref{tab:stats_texto} apresenta as estatísticas descritivas completas.

\begin{table}[!htbp]
	\centering
	\caption{Estatísticas descritivas do comprimento das publicações do corpus}
	\label{tab:stats_texto}
	\begin{tabular}{lrr}
		\hline
		\textbf{Estatística} & \textbf{Caracteres} & \textbf{Tokens} \\
		\hline
		Contagem      & 3.563 & 3.563 \\
		\hline
		Média         &   783 &   123 \\
		\hline
		Desvio padrão &   254 &    40 \\
		\hline
		Mínimo        &   282 &    38 \\
		\hline
		1º quartil    &   596 &    94 \\
		\hline
		Mediana       &   778 &   121 \\
		\hline
		3º quartil    &   950 &   148 \\
		\hline
		Máximo        & 2.503 &   381 \\
		\hline
	\end{tabular}
	\legend{Fonte: Elaborado pelo autor com base nos dados coletados.}
\end{table}

Notavelmente, \textbf{nenhuma publicação} apresenta comprimento inferior a 50 caracteres ou a 5 tokens — limiares usualmente adotados para identificar textos curtos demais para classificação de sentimento. Esse resultado é diretamente explicado pela natureza da fonte: as publicações dos veículos monitorados consistem em resumos e chamadas de notícias que, mesmo no formato de \textit{tweet}, apresentam densidade textual elevada. As distribuições por número de caracteres e por número de tokens são apresentadas na \autoref{fig:histogramas_comprimento}.

\begin{figure}[!htb]
	\caption{\label{fig:histogramas_comprimento}Distribuição do comprimento das publicações por número de caracteres (esquerda) e por número de tokens (direita). As linhas tracejadas vermelha e laranja indicam, respectivamente, os limiares de textos curtos e longos.}
	\begin{minipage}[t]{0.48\textwidth}
		\centering
		\includegraphics[width=\textwidth]{USPSC-TA-Textual/USPSC-Cap3-Metodologia/figuras/Numero_Caracteres.png}
	\end{minipage}
	\hfill
	\begin{minipage}[t]{0.48\textwidth}
		\centering
		\includegraphics[width=\textwidth]{USPSC-TA-Textual/USPSC-Cap3-Metodologia/figuras/Numero_Tokens.png}
	\end{minipage}
	\legend{Fonte: Elaborado pelo autor.}
\end{figure}

No extremo oposto, 678 publicações (19,0\%) ultrapassam 1.000 caracteres e 835 (23,4\%) excedem 150 tokens. Esses textos mais longos não foram descartados, mas devem ser observados com atenção na etapa de pré-processamento: o modelo FinBERT-PT-BR, baseado na arquitetura BERT, aceita sequências de até 512 \textit{tokens} após tokenização com \textit{WordPiece}. Como a tokenização \textit{WordPiece} é mais granular do que a contagem ingênua de palavras por espaços utilizada na AED, o número efetivo de \textit{tokens} BERT é tipicamente superior ao contado aqui. Publicações que excedam o limite do modelo após tokenização serão truncadas automaticamente, e esse impacto constitui uma restrição reconhecida do modelo para textos mais longos.

\subsection{Detecção e Tratamento de Duplicatas}
\label{subsec:duplicatas}

A verificação de duplicatas revelou que \textbf{nenhum registro} apresenta \texttt{tweet\_id} repetido, o que confirma a eficácia da restrição \texttt{UNIQUE} definida no esquema do banco de dados para prevenir a inserção duplicada de publicações durante ciclos de coleta sobrepostos.

A inspeção por conteúdo textual, contudo, identificou \textbf{31 textos com conteúdo idêntico}, envolvendo 54 registros agrupados em 23 grupos — todos provenientes do InfoMoney. Esse padrão é consistente com a prática editorial de republicar a mesma manchete em momentos distintos para ampliar o alcance da publicação, fenômeno documentado na literatura sobre comportamento de contas institucionais em plataformas de mídia social \cite{santana2024cap34}. Para fins da análise de sentimento, os registros duplicados por conteúdo foram mantidos no corpus, uma vez que correspondem a eventos de publicação reais e independentes; sua eventual influência sobre o índice de sentimento agregado é mitigada pela granularidade temporal da agregação diária.

% ============================================================
%  USPSC-Cap3-Metodologia_3.5-Preprocessamento_Textual.tex
%  Seção 3.5 — Pré-processamento Textual
% ============================================================

\section{Pré-processamento Textual}
\label{sec:preprocessamento}

O pré-processamento textual constitui uma etapa crítica em qualquer \textit{pipeline} de PLN, com impacto direto sobre a qualidade das representações vetoriais e, consequentemente, sobre o desempenho do classificador de sentimento \cite{caseli2024cap1}. Para textos provenientes de mídias sociais, essa etapa é particularmente relevante dado o caráter informal, ruidoso e semanticamente denso desses dados, que inclui elementos típicos da linguagem de plataformas digitais --- \textit{hashtags}, menções, URLs e emojis --- que não fazem parte do vocabulário para o qual os modelos de linguagem convencionais foram projetados \cite{santana2024cap34}.

Diferentemente de uma abordagem que aplicaria um único fluxo de limpeza uniforme a todas as abordagens de classificação, este trabalho adota uma \textbf{configuração de pré-processamento específica para cada classificador}, selecionada empiricamente: para cada abordagem, testou-se um conjunto de combinações de etapas de limpeza sobre uma amostra do corpus, comparando a classificação resultante com o rótulo humano e retendo a configuração que produziu a maior taxa de concordância humano-modelo. Essa escolha reconhece que cada abordagem --- léxica ou baseada em modelo de linguagem pré-treinado --- responde de forma distinta ao ruído textual típico do X/Twitter, de modo que a configuração ótima não é necessariamente a mesma entre elas.

As funções de limpeza residem em \textbf{\texttt{app/shared/text\_cleaner.py}}, módulo transversal compartilhado entre os estágios de pré-processamento e classificação, garantindo consistência e evitando duplicação de código. As subseções a seguir descrevem a configuração vencedora de cada abordagem, na ordem em que as etapas são aplicadas.

\subsection{Fluxo de Pré-processamento --- OpLexicon}
\label{subsec:fluxo_oplexicon}

A melhor configuração para o OpLexicon, implementada no método \texttt{preprocess()} da classe \texttt{OpLexiconAnalyzer} (\texttt{app/core/classification/lexicon/op\_lexicon.py}), alcançou \textbf{25,7\% de concordância} humano-modelo e compreende nove etapas sequenciais:

\begin{enumerate}
	\item \textbf{Remoção de URLs}: toda sequência de caracteres iniciada por \texttt{http} ou \texttt{https} é removida do texto, sem substituição por marcador equivalente;
	
	\item \textbf{Remoção de emojis}: os emojis são removidos do texto por meio de \texttt{emoji.replace\_emoji()}, sem substituição por marcador textual ou \textit{name code} equivalente;
	
	\item \textbf{Remoção de menções}: referências a usuários (\texttt{@usuário}) são removidas do texto;
	
	\item \textbf{Remoção do símbolo de \textit{hashtag}}: o caractere \texttt{\#} é removido, mas o termo associado é preservado (e.g., \texttt{\#IBOV} torna-se \texttt{IBOV}) \cite{ranco2015};
	
	\item \textbf{Remoção de pontuação de borda}: caracteres de pontuação de sentença (\texttt{. , ! ?}) são removidos apenas das extremidades de cada \textit{token}, preservando pontuação interna com valor semântico ou estrutural;
	
	\item \textbf{Normalização de espaçamento}: sequências de espaços múltiplos, tabulações e quebras de linha são colapsadas em um único espaço;
	
	\item \textbf{Normalização de caixa com preservação de entidades}: o texto é convertido para letras minúsculas, com exceção de \textit{tickers} de ações (padrão \texttt{[A-Z]\{4,5\}[0-9]\{1,2\}}) e entidades do conjunto \texttt{FINANCIAL\_ENTITIES} (BCB, CVM, IBOV, SELIC, B3, entre outras), que permanecem em maiúsculas;
	
	\item \textbf{Remoção de \textit{stopwords}}: palavras funcionais de alta frequência sem valor semântico para a tarefa são removidas, utilizando a lista de \textit{stopwords} do português da biblioteca \texttt{NLTK} \cite{jurafsky2026};
	
	\item \textbf{Lematização}: as palavras são reduzidas à sua forma canônica (lema) por meio do modelo de língua portuguesa \texttt{pt\_core\_news\_lg} da biblioteca \texttt{spaCy}\footnote{Disponível em: \url{https://spacy.io/}. Acesso em: 12 fev. 2026.}, preservando distinções lexicais relevantes ao domínio (e.g., ``alta'' como substantivo versus ``alto'' como adjetivo). Essa etapa é necessária porque o OpLexicon v3.0 indexa suas entradas por lema, e não por forma flexionada.
\end{enumerate}

\subsection{Fluxo de Pré-processamento --- SentiLex-PT}
\label{subsec:fluxo_sentilex}

A melhor configuração para o SentiLex-PT, implementada no método \texttt{preprocess()} da classe \texttt{SentiLexAnalyzer} (\texttt{app/core/classification/lexicon/senti\_lex.py}), alcançou \textbf{29,5\% de concordância} humano-modelo. Aplica as mesmas oito primeiras etapas descritas para o OpLexicon (\autoref{subsec:fluxo_oplexicon}), na mesma ordem --- remoção de URLs, emojis, menções, símbolo de \textit{hashtag}, pontuação de borda, normalização de espaçamento e de caixa, e remoção de \textit{stopwords} --- \textbf{sem a etapa final de lematização}. Diferentemente do OpLexicon, o SentiLex-PT02 indexa suas entradas tanto por lema quanto por forma flexionada, o que torna a lematização desnecessária para essa abordagem.

\subsection{Fluxo de Pré-processamento --- FinBERT-PT-BR}
\label{subsec:fluxo_finbert}

A configuração vencedora para o FinBERT-PT-BR, implementada no método \texttt{preprocess()} da classe \texttt{FinBertPTBRAnalyzer} (\texttt{app/core/classification/bert/fin bert\_ptbr.py}), alcançou \textbf{40,5\% de concordância} humano-modelo --- a maior entre as abordagens que aplicam alguma etapa de limpeza --- com uma única etapa: \textbf{normalização de espaçamento}, que colapsa sequências de espaços múltiplos, tabulações e quebras de linha em um único espaço.

Esse resultado é notável porque contraria a expectativa inicial de que URLs, emojis, menções e \textit{hashtags} deveriam ser tratados explicitamente antes da inferência. Uma hipótese explicativa é que o tokenizador \textit{WordPiece} interno ao BERT \cite{devlin2019} já é capaz de segmentar esses elementos em subunidades processáveis sem limpeza prévia, e que etapas adicionais de normalização podem remover sinais textuais que, empiricamente, contribuem para a concordância com o julgamento humano nesse corpus específico.

\subsection{Fluxo de Pré-processamento --- BERTimbau Ajustado}
\label{subsec:fluxo_bertimbau}

A configuração vencedora para o BERTimbau ajustado, implementada no método \texttt{preprocess()} da classe \texttt{BERTimbauAnalyzer} (\texttt{app/core/classification/bert/bert\_timbau.py}), alcançou \textbf{89,5\% de concordância} humano-modelo --- a maior entre todas as quatro abordagens --- e compreende sete etapas sequenciais:

\begin{enumerate}
	\item \textbf{Substituição de URLs}: toda sequência iniciada por \texttt{http} ou \texttt{https} é substituída pelo \textit{token} \texttt{[URL]}, preservando o sinal de que um \textit{link} estava presente;
	
	\item \textbf{Substituição de emojis por códigos}: os emojis são convertidos para seus \textit{name codes} por meio de \texttt{emoji.demojize()} (e.g., \includegraphics[width=0.5cm]{USPSC-TA-Textual/png/rotating_light.png} torna-se \texttt{:rotating\_light:}), preservando seu potencial semântico \cite{sprenger2014} de forma processável pelo tokenizador do modelo;
	
	\item \textbf{Substituição de menções}: referências a usuários (\texttt{@usuário}) são substituídas pelo \textit{token} \texttt{[MENTION]};
	
	\item \textbf{Remoção do símbolo de \textit{hashtag}}: o caractere \texttt{\#} é removido, mas o termo associado é preservado;
	
	\item \textbf{Remoção de pontuação de borda}: caracteres de pontuação de sentença são removidos das extremidades de cada \textit{token}, preservando pontuação interna com valor semântico;
	
	\item \textbf{Normalização de espaçamento}: sequências de espaços múltiplos, tabulações e quebras de linha são colapsadas em um único espaço;
	
	\item \textbf{Normalização de caixa com preservação de entidades}: o texto é convertido para minúsculas, com exceção de \textit{tickers} e entidades financeiras, que permanecem em maiúsculas --- consistente com a variante \texttt{cased} do modelo base (\texttt{neuralmind/bert-base-portuguese-cased}), que preserva a distinção entre maiúsculas e minúsculas durante a tokenização \cite{souza2020}.
\end{enumerate}


\subsection{Comparativo entre os Fluxos}
\label{subsec:comparativo_preproc}

O \autoref{tab:comparativo_preproc} sintetiza as diferenças entre as quatro configurações.

\begin{table}[!htbp]
	\caption{Comparativo entre as configurações de pré-processamento vencedoras de cada abordagem}
	\label{tab:comparativo_preproc}
	\begin{tabular}{p{2.6cm}p{2.9cm}p{2.9cm}p{2.3cm}p{2.9cm}}
		\hline
		\textbf{Etapa} & \textbf{OpLexicon} & \textbf{SentiLex-PT} & \textbf{FinBERT-PT-BR} & \textbf{BERTimbau} \\
		\hline
		URLs & Remove & Remove & --- & Substitui por \texttt{[URL]} \\
		\hline
		Emojis & Remove & Remove & --- & Converte p/ \textit{name code} \\
		\hline
		Menções & Remove & Remove & --- & Substitui por \texttt{[MENTION]} \\
		\hline
		Símbolo \texttt{\#} & Remove, mantém termo & Remove, mantém termo & --- & Remove, mantém termo \\
		\hline
		Pontuação de borda & Sim & Sim & --- & Sim \\
		\hline
		Normalização de espaços & Sim & Sim & Sim & Sim \\
		\hline
		Normalização de caixa & Sim, preserva entidades & Sim, preserva entidades & --- & Sim, preserva entidades \\
		\hline
		\textit{Stopwords} & Remove & Remove & --- & --- \\
		\hline
		Lematização & Sim & --- & --- & --- \\
		\hline
		\textbf{Concordância} & \textbf{25,7\%} & \textbf{29,5\%} & \textbf{40,5\%} & \textbf{89,5\%} \\
		\hline
	\end{tabular}
	\legend{Fonte: Elaborado pelo autor com base no código-fonte de \texttt{app/shared/text\_cleaner.py} e dos métodos \texttt{preprocess()} de \texttt{app/core/classification/lexicon/op\_lexicon.py}, \texttt{app/core/classification/lexicon/senti\_lex.py}, \texttt{app/core/classification/bert /finbert\_ptbr.py} e \texttt{app/core/classification/bert/bert\_timbau.py}.}
\end{table}	
	
	
\subsection{Testes Unitários}
\label{subsec:testes_preprocessamento}

O módulo de pré-processamento é acompanhado de um arquivo de testes unitários co-localizado em \texttt{app/shared/text\_cleaner\_tests.py}, executado via \texttt{pytest} a partir da raiz do projeto. Os testes verificam o comportamento de cada função individualmente — cobrindo URLs, emojis, menções, \textit{hashtags}, espaçamento, normalização de caixa, \textit{stopwords} e lematização —, assegurando que alterações futuras no módulo não introduzam regressões silenciosas no fluxo de
pré-processamento.
% ============================================================
%  USPSC-Cap3-Metodologia_3.6-Abordagens_de_Classificacao.tex
%  Seção 3.6 — Abordagens de Classificação de Sentimento
% ============================================================

\section{Abordagens de Classificação de Sentimento}
\label{sec:abordagens}

Em consonância com o Objetivo Específico~4 estabelecido na \autoref{sec:objetivos} — avaliar e comparar o desempenho de diferentes abordagens de classificação de sentimento no domínio financeiro em português —, este trabalho avalia três paradigmas metodológicos distintos: a abordagem \textbf{baseada em léxico}, a abordagem \textbf{baseada em modelo de linguagem pré-treinado especializado por domínio} (FinBERT-PT-BR) e a abordagem \textbf{baseada em modelo de linguagem pré-treinado genérico ajustado por \textit{fine-tuning} em dados locais} (BERTimbau). A distinção entre os dois paradigmas neurais é deliberada: enquanto o FinBERT-PT-BR obtém sua especialização por meio de pré-treinamento sobre um corpus financeiro de larga escala em etapa anterior a este trabalho, o BERTimbau parte de um modelo genérico e adquire especialização \textit{a posteriori}, por meio de ajuste fino sobre o próprio corpus de \textit{tweets} financeiros coletados. Essa dicotomia --- especialização prévia de domínio versus especialização \textit{post hoc} de gênero textual e subdomínio --- constitui o eixo central da comparação empírica apresentada no \autoref{Cap:Ava_Exp}. A inclusão de métodos léxicos permite situar ambas as abordagens neurais em relação às ferramentas de referência disponíveis para o português, orientando recomendações práticas para trabalhos futuros \cite{kearney2014,sousa2025}.


\subsection{Abordagem Baseada em Léxico: OpLexicon}
\label{subsec:abordagem_lexica}

Métodos baseados em léxico (\textit{lexicon-based}) classificam o sentimento de um texto a partir da correspondência de seus termos a um dicionário de polaridade previamente construído, no qual cada palavra ou expressão recebe um escore de valência positiva, negativa ou neutra \cite{liu2020survey}. A principal vantagem dessa abordagem reside na transparência e na ausência de necessidade de dados de treinamento rotulados. Sua limitação central está na incapacidade de capturar negação,
ironia e vocabulário específico do domínio não contemplado pelo léxico \cite{loughran2011,kearney2014}.

\subsubsection{Escolha do Recurso Léxico}
\label{subsubsec:escolha_lexico}

O \textbf{OpLexicon v3.0} foi selecionado como léxico principal por duas razões: (i)~abrangência léxica consideravelmente maior (32.191 entradas de \textit{tokens} não lematizados, cobrindo formas coloquiais recorrentes em \textit{tweets}); e (ii)~construção orientada ao português brasileiro, com termos do jargão financeiro local. O \textbf{SentiLex-PT} \cite{carvalho2015sentilex}, por sua vez, é implementado em paralelo como segunda abordagem léxica (Seção~\ref{subsec:abordagem_sentilex}),  permitindo avaliar o impacto da anotação manual e da cobertura de formas flexionadas sobre o desempenho no domínio financeiro.


\subsubsection{Estrutura do Léxico}
\label{subsubsec:estrutura_lexico}

O OpLexicon v3.0 é distribuído como arquivo de texto tabular (\textit{tab-separated values}) sem cabeçalho, com quatro colunas: \texttt{term} (lema), \texttt{type} (classe gramatical), \texttt{polarity} (polaridade inteira: $+1$, $-1$ ou $0$) e \texttt{polarity\_revision} (origem da entrada: \texttt{A} para revisão manual ou \texttt{C} para atribuição automática). A indexação por lemas — e não por formas flexionadas — é a razão pela qual a lematização constitui etapa obrigatória do pré-processamento para esta abordagem, conforme detalhado na \autoref{subsec:fluxo_lexico}.

\subsubsection{Implementação e Protocolo de Pontuação}
\label{subsubsec:protocolo_oplexicon}

A implementação foi realizada no quinto estágio do \textit{pipeline} por meio da classe \texttt{OpLexiconAnalyzer}, definida em \texttt{app/core/processing/lexicon/op\_lexicon.py}. A classe herda de \texttt{BaseSentimentAnalyzer} — a mesma classe abstrata (ABC) que fundamenta o
\texttt{FinBertPTBRAnalyzer} —, garantindo interface uniforme entre as duas abordagens e extensibilidade para classificadores futuros.

O protocolo de pontuação compreende as seguintes etapas:

\begin{enumerate}
	\item \textbf{Seleção dos registros}: o método \texttt{run()} recupera exclusivamente as publicações que já possuem classificação humana em \texttt{tweets\_classification} mas ainda não foram processadas pelo OpLexicon, garantindo que a comparação opere sobre o mesmo
	\textit{gold standard} utilizado pelo FinBERT-PT-BR;
	
	\item \textbf{Pré-processamento}: cada texto é submetido ao método \texttt{preprocess()}, que aplica o fluxo completo descrito na \autoref{subsec:fluxo_lexico}, incluindo lematização via \texttt{spaCy} — etapa obrigatória para maximizar a taxa de correspondência com as entradas do léxico, que indexa lemas e não formas flexionadas;
	
	\item \textbf{Carregamento do léxico}: o método \texttt{load\_model()} lê o arquivo \texttt{lexico\_v3.0.txt} e constrói um dicionário em memória no formato
	\texttt{\{lema: polaridade\}}, com tempo de carga desprezível e sem demanda de GPU;
	
	\item \textbf{Pontuação agregada}: para cada \textit{token} do texto pré-processado, consulta-se o dicionário léxico. \textit{Tokens} ausentes recebem pontuação zero e são excluídos da agregação. A polaridade final é calculada pela média aritmética dos escores dos \textit{tokens} com entrada no léxico, conforme a \autoref{eq:oplexicon_score}:
\end{enumerate}

\begin{equation}
	\label{eq:oplexicon_score}
	\bar{p} \;=\; \frac{1}{|T^{*}|} \sum_{t \,\in\, T^{*}} \mathrm{pol}(t)
\end{equation}

\noindent onde $T^{*} = \{t \in T \mid \mathrm{pol}(t) \neq 0\}$ é o subconjunto de \textit{tokens} com polaridade não-nula e $\mathrm{pol}(t)$ é a polaridade do \textit{token} $t$
no léxico. Caso $T^{*} = \emptyset$, o texto é classificado como Neutro com escore zero.

\begin{enumerate}
	\setcounter{enumi}{4}
	\item \textbf{Limiarização}: a classe final é determinada por um limiar $\theta = 0{,}05$ aplicado à média $\bar{p}$: se $\bar{p} > \theta$, a publicação é classificada como Positivo; se $\bar{p} < -\theta$, como Negativo; caso contrário, como Neutro. O valor reduzido do limiar reflete a alta proporção esperada de \textit{tokens} financeiros sem correspondência no léxico (siglas de ativos, nomes de emissores, jargão de derivativos), que dilui a magnitude da média
	sem necessariamente indicar neutralidade semântica;
	
	\item \textbf{Persistência}: os resultados são armazenados em \texttt{tweets\_classification} com \texttt{classificator = "OpLexicon"} e o respectivo escore, permitindo comparação direta com as anotações humanas e com o FinBERT-PT-BR na etapa de avaliação (\autoref{subsec:metricas_avaliacao}).
\end{enumerate}

% ------------------------------------------------------------------
\subsection{Abordagem Baseada em Léxico: SentiLex-PT}
\label{subsec:abordagem_sentilex}
% ------------------------------------------------------------------

\subsubsection{Fundamentação da Escolha}
\label{subsubsec:justificativa_sentilex}

O \textbf{SentiLex-PT} \cite{carvalho2015sentilex} foi incorporado como segunda abordagem léxica com o objetivo de ampliar a validade comparativa da avaliação. Enquanto o OpLexicon~v3.0 é construído sobre um processo semiautomático voltado ao português brasileiro, o SentiLex-PT decorre de um esforço de anotação manual, o que lhe confere características complementares relevantes para o domínio financeiro:

\begin{enumerate}
	\item \textbf{Anotação manual}: os 7.014 lemas (expandidos para 82.347 formas flexionadas na versão SentiLex-PT02) foram anotados manualmente por especialistas segundo a taxonomia \textit{Manually Annotated} (\texttt{ANOT:MAN}), resultando em um recurso de alta precisão intrínseca \cite{carvalho2015sentilex};

	\item \textbf{Cobertura de formas flexionadas}: a versão SentiLex-PT02, adotada neste trabalho, contém formas nominais, adjetivais e verbais flexionadas, dispensando etapa de lematização e tornando-a mais adequada ao processamento de \textit{tweets} --- textos curtos com vocabulário coloquial e morfologia variada;

	\item \textbf{Validade científica consolidada}: o SentiLex-PT é amplamente utilizado como referência em estudos de análise de sentimento em português \cite{liu2020survey}, permitindo contextualizar os resultados obtidos neste trabalho em relação à literatura.
\end{enumerate}

\subsubsection{Estrutura do Léxico}
\label{subsubsec:estrutura_sentilex}

O SentiLex-PT02 é distribuído como arquivo de texto simples, com uma entrada por linha no formato:

\begin{center}
	\texttt{forma.PoS,POL:\textit{N};ANOT:\textit{TAG}}
\end{center}

\noindent onde \texttt{forma} é a palavra flexionada em letras minúsculas, \texttt{PoS} é a classe gramatical (\texttt{Adj}, \texttt{N}, \texttt{V}, entre outras), \texttt{POL} é a polaridade inteira (\texttt{1}~=~positivo, \texttt{-1}~=~negativo, \texttt{0}~=~neutro) e \texttt{ANOT} é o tipo de anotação. A \autoref{tab:sentilex_exemplos} apresenta exemplos representativos do léxico no domínio financeiro.

\begin{table}[!htbp]
	\caption{Exemplos de entradas do SentiLex-PT02 no domínio financeiro}
	\label{tab:sentilex_exemplos}
	\centering
	\begin{tabular}{lllc}
		\hline
		\textbf{Forma} & \textbf{PoS} & \textbf{ANOT} & \textbf{POL} \\
		\hline
		lucro      & N   & MAN & \phantom{$-$}1 \\
		\hline
		crescimento & N  & MAN & \phantom{$-$}1 \\
		\hline
		valorização & N  & MAN & \phantom{$-$}1 \\
		\hline
		queda      & N   & MAN & $-$1 \\
		\hline
		prejuízo   & N   & MAN & $-$1 \\
		\hline
		crise      & N   & MAN & $-$1 \\
		\hline
		empresa    & N   & MAN & \phantom{$-$}0 \\
		\hline
		mercado    & N   & MAN & \phantom{$-$}0 \\
		\hline
	\end{tabular}
	\fonte{Adaptado de \citeonline{carvalho2015sentilex}.}
\end{table}

\subsubsection{Implementação e Protocolo de Pontuação}
\label{subsubsec:protocolo_sentilex}

A implementação foi realizada no quinto estágio do \textit{pipeline} por meio da classe \texttt{SentiLexAnalyzer}, definida em \texttt{app/core/processing/lexicon/senti\_lex.py}. A classe herda de \texttt{BaseSentimentAnalyzer} (\texttt{app/core/processing/base\_analyzer.py}), garantindo interface uniforme com o \texttt{OpLexiconAnalyzer} e o \texttt{FinBertPTBRAnalyzer}.

O protocolo de pontuação compreende as seguintes etapas:

\begin{enumerate}
	\item \textbf{Seleção dos registros}: o método \texttt{run()} recupera exclusivamente as publicações que já possuem classificação humana em \texttt{tweets\_classification}, utilizadas como \textit{gold standard};

	\item \textbf{Pré-processamento}: cada texto é submetido ao método \texttt{preprocess()}, que aplica o fluxo descrito na \autoref{sec:preprocessamento} --- remoção de URLs, menções, caracteres especiais e conversão para letras minúsculas. Diferentemente do \texttt{OpLexiconAnalyzer}, não é realizada lematização via \texttt{spaCy}, pois o SentiLex-PT02 já contém formas flexionadas;

	\item \textbf{Carregamento do léxico}: o método \texttt{load\_model()} lê o arquivo \texttt{data/SentiLex- PT02.txt} e constrói um dicionário em memória no formato \texttt{\{forma: polaridade\}}, com tempo de carga desprezível e sem demanda de GPU;

	\item \textbf{Pontuação com janela de negação}: para cada \textit{token} do texto pré-processado, consulta-se o dicionário léxico. A fim de capturar negações simples --- padrão frequente em textos financeiros, como \textit{``não houve crescimento''} ---, aplica-se uma janela de negação de três \textit{tokens}: quando um \textit{token} pertencente ao conjunto de marcadores de negação\footnote{Conjunto de marcadores: \textit{não, nunca, jamais, nem, nenhum, nenhuma, tampouco, sequer}.} precede uma palavra com polaridade não-nula dentro desse intervalo, o sinal da polaridade é invertido. A pontuação final é dada por:

	\begin{equation}
		\label{eq:sentilex_score}
		\bar{p} = \sum_{t \in T^{*}} \mathrm{pol\_neg}(t)
	\end{equation}

	\noindent onde $T^{*} = \{t \in T \mid \mathrm{pol}(t) \neq 0\}$ é o subconjunto de \textit{tokens} com polaridade não-nula e $\mathrm{pol\_neg}(t)$ é a polaridade do \textit{token} $t$ após aplicação da janela de negação. Caso $T^{*} = \emptyset$, o texto é classificado como Neutro com escore zero;

	\item \textbf{Limiarização}: a classe final é determinada diretamente pelo sinal da pontuação acumulada $\bar{p}$: se $\bar{p} > 0$, a publicação é classificada como Positivo; se $\bar{p} < 0$, como Negativo; se $\bar{p} = 0$, como Neutro. O escore de confiança reportado é dado por $|\bar{p}| / |T|$, normalizado ao intervalo $[0, 1]$, análogo ao \textit{score} de probabilidade retornado pelo FinBERT-PT-BR;

	\item \textbf{Persistência}: os resultados são armazenados em \texttt{tweets\_classification} com \texttt{classificator = ``SentiLex-PT''} e o respectivo escore, permitindo comparação direta com as demais abordagens conforme descrito na \autoref{subsec:metricas_avaliacao}.
\end{enumerate}


\subsection{Abordagem Baseada em Modelo de Linguagem Pré-treinado: FinBERT-PT-BR}
\label{subsec:finbert}

\subsubsection{Fundamentação da Escolha}
\label{subsubsec:justificativa_finbert}

O \textbf{FinBERT-PT-BR} \cite{santos2023finbert} foi selecionado como ferramenta central da abordagem neural por três razões fundamentais, já antecipadas no \autoref{Cap:Fundamentacao_Teorica}:

\begin{enumerate}
	\item \textbf{Especialização no domínio}: o modelo foi pré-treinado sobre um corpus de 1,4 milhão de textos financeiros em português, incluindo notícias, relatórios e publicações de veículos especializados, o que lhe confere representações contextuais calibradas para o vocabulário técnico do mercado financeiro brasileiro;
	
	\item \textbf{Desempenho superior no estado da arte}: nos experimentos reportados por \citeonline{santos2023finbert}, o FinBERT-PT-BR superou modelos de linguagem genéricos e abordagens léxicas em tarefas de classificação de sentimento financeiro em português;
	
	\item \textbf{Disponibilidade pública}: o modelo está disponível publicamente no repositório \textit{Hugging Face} sob o identificador \texttt{lucas-leme/FinBERT-PT-BR}, garantindo a reprodutibilidade da metodologia proposta.
\end{enumerate}

\subsubsection{Arquitetura do Modelo}
\label{subsubsec:arquitetura_finbert}

O FinBERT-PT-BR é construído sobre a arquitetura \textbf{BERT} (\textit{Bidirectional Encoder Representations from Transformers}) \cite{devlin2019}, que implementa o mecanismo de atenção multi-cabeça proposto por \citeonline{vaswani2017}. A característica central do BERT — a codificação bidirecional do contexto — permite que a
representação de cada \textit{token} incorpore informações tanto dos termos que o precedem quanto dos que o sucedem na sequência, o que é particularmente relevante para a interpretação de negações e modificadores de intensidade frequentes no discurso financeiro (e.g., ``\textit{não} sobe'', ``\textit{queda} acentuada'').

O modelo base utilizado para o \textit{fine-tuning} é o \textbf{BERTimbau} \cite{souza2020}, pré-treinado sobre um corpus de 2,68 bilhões de palavras em português brasileiro. Sobre esse modelo base, \citeonline{santos2023finbert} realizaram o \textit{fine-tuning} para a tarefa de classificação de sentimento financeiro em três classes:
\textbf{Positivo}, \textbf{Negativo} e \textbf{Neutro}.

\subsubsection{Implementação e Protocolo de Inferência}
\label{subsubsec:protocolo_inferencia}

A inferência foi implementada no quinto estágio do \textit{pipeline} (\texttt{app/core/processing /bert/}), por meio da classe \texttt{FinBertPTBRAnalyzer}, definida em \texttt{app/core/processing/ bert/finbert\_ptbr.py}. A classe herda de \texttt{BaseSentimentAnalyzer} (\texttt{app/core/process ing/base\_analyzer.py}), uma classe abstrata (ABC) que estabelece o contrato de interface para qualquer analisador de sentimento no \textit{pipeline}, garantindo extensibilidade para futuros modelos.

O protocolo de inferência compreende as seguintes etapas:

\begin{enumerate}
	\item \textbf{Seleção dos registros}: o método \texttt{run()} recupera via \texttt{app/shared/db/database.py} exclusivamente as publicações que já possuem classificação humana em \texttt{tweets\_classification} mas ainda não foram processadas pelo FinBERT-PT-BR. Esse critério garante que a inferência do modelo opere sobre o mesmo subconjunto utilizado como \textit{gold standard}, permitindo comparação direta entre as duas abordagens;
	
	\item \textbf{Pré-processamento}: cada texto é submetido ao método \texttt{preprocess()} da própria classe, que aplica o fluxo reduzido descrito na \autoref{subsec:fluxo_finbert} — substituição de URLs e menções, remoção de emojis, remoção do símbolo de \textit{hashtag}, normalização de espaços e normalização de caixa com preservação de entidades financeiras;
	
	\item \textbf{Carregamento do modelo}: o tokenizador e o modelo são instanciados via \texttt{AutoToke nizer} e \texttt{BertForSequenceClassification}, encapsulados em um objeto \texttt{pipeline} da biblioteca \texttt{transformers} com \texttt{task="text-classification"};
	
	\item \textbf{Inferência}: cada publicação é classificada pelo método \texttt{predict\_text()}, com \texttt{batch\_size=32}. O modelo retorna o \textit{label} predito e o escore de confiança (\textit{score}) associado;
	
	\item \textbf{Mapeamento de \textit{labels}}: os rótulos em inglês retornados pelo modelo são normalizados para o português: \texttt{POSITIVE} → \texttt{positivo}, 	\texttt{NEGATIVE} → \texttt{negativo}, \texttt{NEUTRAL} → \texttt{neutro};
	
	\item \textbf{Persistência e visualização}: os resultados são armazenados em \texttt{tweets\_classifi cation} com \texttt{classificator = "FinBERT-PT-BR"} e o respectivo 	\texttt{score}. O painel Streamlit de processamento (\texttt{app/dashboard/pages/processing.py}), invocado via \texttt{make dashboard}, apresenta a comparação entre o sentimento predito pelo modelo e o sentimento anotado pelo humano tweet a tweet, além de métricas de concordância.
\end{enumerate}

\subsection{Abordagem Baseada em Modelo de Linguagem Pré-treinado com \textit{Fine-tuning}: BERTimbau}
\label{subsec:bertimbau}

\subsubsection{Fundamentação da Escolha}
\label{subsubsec:justificativa_bertimbau}

O \textbf{BERTimbau} \cite{souza2020} foi selecionado como modelo base para a abordagem de \textit{fine-tuning} local por três razões complementares:

\begin{enumerate}
	\item \textbf{Estado da arte para o português brasileiro genérico}: o BERTimbau é o modelo BERT pré-treinado de referência para a língua portuguesa, treinado sobre um corpus de 2,68 bilhões de palavras extraídas de fontes diversificadas --- Wikipedia em português, Common Crawl e BrWaC \cite{souza2020}. Sua adoção como base de \textit{fine-tuning} é amplamente consolidada na literatura de PLN em português, inclusive como ponto de partida do próprio FinBERT-PT-BR \cite{santos2023finbert};
	
	\item \textbf{Endereçamento direto do desalinhamento de domínio}: o FinBERT-PT-BR foi pré-ajustado sobre notícias e relatórios financeiros, gênero textual estruturalmente distinto das publicações do X/Twitter. O \textit{fine-tuning} do BERTimbau diretamente sobre o corpus de \textit{tweets} coletados --- incluindo suas características estilísticas de informalidade, densidade de abreviações e uso de emojis e \textit{hashtags} --- aborda essa limitação de forma explícita, produzindo um modelo alinhado ao gênero e ao subdomínio específico avaliado;
	
	\item \textbf{Adequação ao tamanho do corpus}: a versão \texttt{base} do BERTimbau (110 milhões de parâmetros) foi preferida à versão \texttt{large} (335 milhões de parâmetros) pela maior adequação ao tamanho do corpus rotulado disponível. Modelos de maior capacidade requerem volumes de dados proporcionalmente maiores para generalização estável, e o uso da versão \texttt{large} sobre um corpus de algumas centenas ou poucos milhares de exemplos aumentaria o risco de \textit{overfitting} sem ganho garantido de desempenho \cite{devlin2019}.
\end{enumerate}

\subsubsection{Arquitetura e Protocolo de \textit{Fine-tuning}}
\label{subsubsec:arquitetura_bertimbau}

O modelo base utilizado é o \texttt{neuralmind/bert-base-portuguese-cased}, disponível publicamente no repositório \textit{Hugging Face}. Sobre esse modelo, realiza-se o \textit{fine-tuning} para a tarefa de classificação de sequência em três classes --- \textbf{Positivo}, \textbf{Neutro} e \textbf{Negativo} ---, por meio da adição de uma camada de classificação linear sobre a representação do \textit{token} especial \texttt{[CLS]}, que codifica o sentido global da sequência de entrada \cite{devlin2019}.

O protocolo de \textit{fine-tuning} é implementado no script \texttt{app/core/processing/bert /bert\_timbau\_fine\_tuner.py} e compreende as seguintes decisões de configuração:

\begin{itemize}

	\item \textbf{Particionamento}: o corpus rotulado é particionado em etapa dedicada, anterior ao \textit{fine-tuning}, conforme o protocolo descrito na \autoref{subsec:particionamento}. O particionamento é gerenciado pela classe \texttt{DatasetSplitRepository} (\texttt{app/shared/db/dataset\_split.py}), que persiste o resultado na tabela \texttt{dataset\_split} em dois conjuntos mutuamente exclusivos: (i)~\textbf{hold-out de teste} ($\approx$~15\% dos registros), extraído por amostragem estratificada com \texttt{random\_state=42} e reservado exclusivamente para a avaliação comparativa final descrita na \autoref{subsec:metricas_avaliacao}, marcado com \texttt{split=`test'} e \texttt{fold=NULL}; e (ii)~\textbf{conjunto de treinamento} ($\approx$~85\%), subdividido em $K=4$ folds por validação cruzada estratificada (\texttt{StratifiedKFold}, $K=4$, \texttt{shuffle=True}, \texttt{random\_state=42}), marcado com \texttt{split=`train'} e \texttt{fold}~$\in \{1, 2, 3, 4\}$;
	
	\item \textbf{Validação cruzada e seleção do modelo}: o \textit{fine-tuning} é executado $K=4$ vezes, uma por fold. Em cada iteração~$k$, os registros com \texttt{fold}~$\neq k$ ($\approx$~64\% do total rotulado) formam o conjunto de treinamento efetivo, e os registros com \texttt{fold}~$= k$ ($\approx$~21\%) formam o conjunto de validação. Ao término das $K$ iterações, o modelo correspondente ao fold com maior F1-Score macro na validação é selecionado como modelo definitivo e persistido em \texttt{models/bert-timbau-sentiment/}, acompanhado do relatório de classificação e do resumo do K-Fold em \texttt{models/bert-timbau-sentiment/eval/};
	
	\item \textbf{Ponderação de classes}: dado o desbalanceamento esperado em favor da classe neutra, os pesos por classe são calculados via \texttt{compute\_class\_weight(class\_weight=`balanced')} da biblioteca \texttt{scikit-learn} e injetados na função de perda (\texttt{CrossEntropyLoss}) por meio de uma subclasse \texttt{WeightedTrainer}, garantindo que o modelo não colapse para predizer sistematicamente a classe majoritária;
	
	\item \textbf{Hiperparâmetros}: taxa de aprendizado $5 \times 10^{-5}$, tamanho de lote (\textit{batch size}) 8, comprimento máximo de sequência 512 \textit{tokens}, razão de \textit{warmup} 0,1 e decaimento de pesos (\textit{weight decay}) $10^{-2}$. Esses valores foram determinados por busca empírica sobre o conjunto de validação cruzada e são consistentes com as recomendações dos autores do BERT para \textit{fine-tuning} em conjuntos de dados de tamanho moderado \cite{devlin2019};
	
	\item \textbf{\textit{Early stopping}}: o treinamento de cada fold é limitado a doze épocas, com critério de parada antecipada (\textit{patience}~=~2) monitorando o F1-Score macro no conjunto de validação do fold corrente. O \textit{checkpoint} com melhor F1-Score macro é restaurado automaticamente ao término de cada iteração, prevenindo o \textit{overfitting} à distribuição de treinamento;
	
	\item \textbf{Reprodutibilidade}: a semente aleatória \texttt{random\_state=42} é fixada em todas as operações estocásticas do treinamento (\texttt{seed} do \texttt{TrainingArguments}), garantindo que execuções sucessivas sobre o mesmo corpus produzam resultados idênticos.
\end{itemize}

O modelo ajustado, o tokenizador e o mapeamento de rótulos (\texttt{id2label.json}) são persistidos em \texttt{models/bert-timbau-sentiment/}, e o relatório de classificação e a matriz de confusão no conjunto de validação são salvos em \texttt{models/bert-timbau-sentiment/eval/}, para referência na análise apresentada no \autoref{Cap:Ava_Exp}.

\subsubsection{Implementação e Protocolo de Inferência}
\label{subsubsec:protocolo_inferencia_bertimbau}

A inferência é implementada pela classe \texttt{BERTimbauAnalyzer}, definida em \texttt{app/core/ processing/bert/bert\_timbau.py}. Assim como \texttt{FinBertPTBRAnalyzer}, a classe herda de \texttt{BaseSentimentAnalyzer} (\texttt{app/core/processing/base\_analyzer.py}), garantindo que ambos os analisadores compartilhem o mesmo contrato de interface e possam ser alternados pelo painel de processamento sem modificações adicionais. O protocolo de inferência é análogo ao descrito na \autoref{subsubsec:protocolo_inferencia}, com as seguintes particularidades:

\begin{enumerate}
	\item \textbf{Pré-requisito de treinamento}: a classe verifica na inicialização a existência do diretório \texttt{models/bert-timbau-sentiment/}; caso ausente, emite um erro orientando o usuário a executar \texttt{python -m pipeline.05\_processing.bert\_timbau\_fine\_tuner} (atalho: \texttt{make finetune}) antes da inferência;
	
	\item \textbf{Carregamento do modelo}: tokenizador e modelo são carregados diretamente do diretório local via \texttt{AutoTokenizer} e \texttt{AutoModelForSequenceClassification}, encapsulados em objeto \texttt{pipeline} da biblioteca \texttt{transformers} com \texttt{task="text-classifi cation"};
	
	\item \textbf{Mapeamento de rótulos}: o modelo fine-tuned emite rótulos em português (\texttt{positivo}, \texttt{negativo}, \texttt{neutro}) diretamente via \texttt{id2label}, dispensando a etapa de normalização de inglês para português aplicada ao FinBERT-PT-BR;
	
	\item \textbf{Persistência}: os resultados são armazenados em \texttt{tweets\_classification} com \texttt{classificator = "BERTimbau"} e o respectivo escore de confiança, seguindo o mesmo esquema de persistência descrito na \autoref{subsubsec:protocolo_inferencia} para o FinBERT-PT-BR.
\end{enumerate}


\subsection{Métricas de Avaliação das Abordagens}
\label{subsec:metricas_avaliacao}

A comparação entre as quatro abordagens --- OpLexicon, SentiLex-PT, FinBERT-PT-BR e BERTimbau ajustado --- será realizada exclusivamente sobre o conjunto de teste \textit{hold-out} definido na \autoref{subsec:particionamento}, que constitui o \textit{gold standard} efetivo da avaliação. Essa restrição é necessária para garantir comparabilidade: como o BERTimbau ajustado viu parte do corpus rotulado durante o treinamento, avaliá-lo sobre o conjunto completo produziria estimativas otimistas de desempenho incomparáveis às das demais abordagens. A implementação das métricas está organizada no sexto estágio do \textit{pipeline} (\texttt{app/core/evaluation/metrics.py}), invocado via \texttt{make evaluate}. As métricas calculadas são:

\begin{itemize}
	\item \textbf{Acurácia}: proporção de classificações corretas sobre o total de instâncias do conjunto de teste;
	
	\item \textbf{Precisão, Revocação e F1-Score} por classe: métricas calculadas individualmente para as classes Positivo, Negativo e Neutro, permitindo identificar vieses de classificação em classes específicas;
	
	\item \textbf{F1-Score macro}: média aritmética não ponderada dos F1-Scores de cada classe, tratando as classes como igualmente importantes independentemente de seu tamanho --- métrica de referência para comparação entre abordagens em corpora desbalanceados;
	
	\item \textbf{Matriz de confusão}: representação tabular da distribuição de acertos e erros por classe, útil para identificar padrões sistemáticos de classificação incorreta entre as três abordagens.
\end{itemize}

A limitação inerente à anotação por avaliador único --- ausência de cálculo de concordância interanotadores (\textit{inter-annotator agreement}) --- é reconhecida como uma restrição metodológica do presente trabalho.
% ============================================================
%  USPSC-Cap3-Metodologia_3.7-Pipeline_Integrado.tex
%  Seção 3.7 — Pipeline Integrado de PLN
% ============================================================

\section{\textit{Pipeline} Integrado de PLN}
\label{sec:pipeline}

O conjunto das etapas descritas nas seções anteriores foi organizado em um \textit{pipeline} integrado de PLN de ponta a ponta (\textit{end-to-end}), implementado em Python. A modularidade da implementação garante que cada etapa possa ser executada, auditada e substituída de forma independente, o que é um requisito central para a reprodutibilidade da metodologia proposta \cite{santos2023finbert}.

\subsection{Visão Geral do \textit{Pipeline}}
\label{subsec:visao_pipeline}

O \textit{pipeline} é estruturado em seis etapas sequenciais, conforme ilustrado na \autoref{fig:pipeline}:

\begin{enumerate}
	\item \textbf{Coleta via API}: interação com o \textit{endpoint} \texttt{GET /2/users/\{id\}/tweets} da API v2 do X/Twitter, com autenticação por Bearer Token e paginação iterativa dos resultados, conforme detalhado na \autoref{subsec:coleta_api};
	
	\item \textbf{Persistência em banco de dados}: os registros retornados pela API são armazenados diretamente no banco de dados PostgreSQL — provisionado em contêiner Docker — nas tabelas
	\texttt{tweets} e \texttt{collection\_log}, eliminando a necessidade de arquivos intermediários em disco;
	
	\item \textbf{Rotulagem manual}: as publicações persistidas são exibidas sequencialmente na ferramenta de anotação (\textit{Classifier App}), onde o anotador atribui relevância financeira e polaridade de sentimento a cada registro, conforme o protocolo descrito na \autoref{sec:rotulagem};
	
	\item \textbf{AED e pré-processamento textual}: o corpus rotulado é inspecionado estatisticamente e submetido às transformações linguísticas descritas nas Seções~\ref{sec:eda} 	e~\ref{sec:preprocessamento};
	
	\item \textbf{Classificação de sentimento}: o corpus pré-processado é submetido às quatro abordagens de classificação de sentimento descritas na \autoref{sec:abordagens}: (a)~abordagem léxica via OpLexicon~v3.0; (b)~abordagem léxica via SentiLex-PT; (c)~abordagem neural via FinBERT-PT-BR; e (d)~abordagem neural via BERTimbau ajustado. Todas gravam os resultados em \texttt{tweets\_classification} com o respectivo campo \texttt{classificator}, permitindo comparação direta com o \textit{gold standard} humano na etapa de avaliação;

	\item \textbf{Agregação e análise}: as classificações individuais são agregadas temporalmente e os resultados são comparados entre as abordagens e analisados à luz do contexto macroeconômico do período, conforme apresentado no \autoref{Cap:Ava_Exp}.
\end{enumerate}

\begin{figure}[!htb]
	\centering
	\caption{\label{fig:pipeline}Diagrama do \textit{pipeline} integrado de PLN para análise de sentimento de publicações financeiras no X/Twitter}
	\resizebox{\textwidth}{!}{%
		\begin{tikzpicture}[
			node distance = 1.0cm and 0.5cm,
			box/.style = {rectangle, rounded corners=4pt, draw=black,
				fill=gray!10, text width=2.6cm, align=center,
				minimum height=1.1cm, font=\small},
			arrow/.style = {->, thick, >=stealth}
			]
	
			\node[box] (coleta) {Coleta\\via API v2};
			\node[box, right=of coleta] (db) {Persistência\\PostgreSQL};
			\node[box, right=of db] (rotul) {Rotulagem\\Manual};
			\node[box, right=of rotul] (preproc) {AED e\\Pré-proc.};
			\node[box, right=of preproc] (classif) {Classificação\\Sentimento};
			\node[box, right=of classif] (agreg) {Agregação\\e Análise};
			
			\node[above=0.35cm of coleta, font=\scriptsize] {API X/Twitter};
			\node[above=0.35cm of agreg,  font=\scriptsize] {Índice de Sentimento};
			
			\draw[arrow] (coleta)   -- (db);
			\draw[arrow] (db)       -- (rotul);
			\draw[arrow] (rotul)    -- (preproc);
			\draw[arrow] (preproc)  -- (classif);
			\draw[arrow] (classif)  -- (agreg);
			
			\node[below=0.30cm of coleta,  font=\scriptsize] {Bearer Token};
			\node[below=0.30cm of db,      font=\scriptsize] {tabela \texttt{tweets}};
			\node[below=0.30cm of rotul,   font=\scriptsize] {rótulos humanos};
			\node[below=0.30cm of preproc, font=\scriptsize] {corpus norm.};
			\node[below=0.30cm of classif, font=\scriptsize] {prob. classes};
		\end{tikzpicture}%
	}
	\legend{Fonte: Elaborado pelo autor.}
\end{figure}

\subsection{Ambiente de Execução e Ferramentas}
\label{subsec:ambiente_execucao}

A implementação foi realizada na linguagem \textbf{Python 3.11}, escolhida por sua ampla adoção na comunidade de PLN e aprendizado de máquina e pela disponibilidade de ecossistema maduro de bibliotecas especializadas. O \autoref{tab:bibliotecas} lista as principais dependências utilizadas e suas respectivas funções no \textit{pipeline}.

\begin{table}[!htb]
	\caption[Principais bibliotecas Python utilizadas no \textit{pipeline} e respectivas funções]{Principais bibliotecas Python utilizadas no \textit{pipeline} e respectivas funções}
	\label{tab:bibliotecas}
	\centering
	\begin{tabular}{ccp{10cm}}
		\hline
		\textbf{Biblioteca} & \textbf{Versão} & \textbf{Função no \textit{Pipeline}} \\
		
		\hline
		\texttt{requests} & $\geq$ 2.31 & Requisições HTTP à API v2 do X/Twitter com Bearer Token \\
		
		\hline
		\texttt{psycopg2} & $\geq$ 2.9 & Interface de comunicação com o banco de dados PostgreSQL \\
		
		\hline
		\texttt{streamlit} & $\geq$ 1.35 & Ferramenta de anotação manual (\textit{Classifier App}) e painel de AED (\texttt{eda.py}) \\
		
		\hline
		\texttt{pandas} & $\geq$ 2.0 & Manipulação tabular do corpus, AED e construção do índice temporal \\
		
		\hline
		\texttt{transformers} & $\geq$ 4.40 & Carregamento e inferência com o modelo FinBERT-PT-BR \\
		
		\hline
		\texttt{torch} & $\geq$ 2.2 & \textit{Backend} de inferência neural (PyTorch) \\
		
		\hline
		\texttt{scikit-learn} & $\geq$ 1.4 & Particionamento estratificado do corpus, cálculo de métricas de avaliação e pesos por classe para \textit{fine-tuning} \\
		
		\hline
		\texttt{accelerate} & $\geq$ 0.28 & \textit{Backend} de paralelização do \texttt{Trainer} da biblioteca \texttt{transformers}, requerido pelo protocolo de \textit{fine-tuning} do BERTimbau \\
		
		\hline
		\texttt{datasets} & $\geq$ 2.18 & Estruturação eficiente de conjuntos de dados para o ciclo de \textit{fine-tuning} supervisionado \\
		
		\hline
		\texttt{nltk} & $\geq$ 3.8 & Remoção de \textit{stopwords} para abordagem léxica \\
		
		\hline
		\texttt{spacy} & $\geq$ 3.7 & Lematização com modelo \texttt{pt\_core\_news\_lg} \\
		
		\hline
		\texttt{emoji} & $\geq$ 2.10 & Transliteração de emojis para texto \\
		
		\hline
		\texttt{plotly} & $\geq$ 5.20 & Visualizações interativas da AED e do índice de sentimento \\
		\hline
	\end{tabular}
	\legend{Fonte: Elaborado pelo autor.}
\end{table}

O código-fonte do \textit{pipeline} está organizado como um pacote Python versionado, com dependências declaradas em arquivo \texttt{requirements.txt}, de modo a facilitar a reprodução dos experimentos em ambientes distintos.
%\input{USPSC-TA-Textual/USPSC-Cap3-Metodologia/USPSC-Cap3-Metodologia_3.8-Indice_de_Sentimento.tex}
% ============================================================
%  USPSC-Cap3-Metodologia_3.9-Limitacoes.tex
%  Seção 3.9 — Limitações da Pesquisa
% ============================================================

\section{Limitações da Pesquisa}
\label{sec:limitacoes}

A identificação e discussão honesta das limitações metodológicas integra os requisitos de rigor científico do método \cite{caseli2024cap1}. As limitações do presente trabalho operam em quatro dimensões: fonte de dados, representatividade do corpus, restrições do modelo e escopo analítico.

\subsection{Limitações Relativas à Fonte e ao Período de Dados}
\label{subsec:lim_fonte}

\textbf{Janela temporal restrita.} O corpus cobre aproximadamente quatro meses de dados (outubro de 2025 a fevereiro de 2026), período que, embora suficiente para os objetivos descritivos desta pesquisa, impede generalizações sobre o comportamento do sentimento ao longo de ciclos econômicos mais longos. A limitação decorre das restrições de custo de acesso à API histórica do X/Twitter, conforme detalhado na \autoref{subsec:periodo}.

\textbf{Dependência da API do X/Twitter.} A coleta de dados via API está sujeita a alterações de política e estrutura de preços da plataforma, o que representa um risco de reprodutibilidade para pesquisas futuras que tentem replicar o protocolo com o mesmo instrumento de coleta. Mudanças nas políticas de acesso do X/Twitter têm sido frequentes desde a aquisição da plataforma em 2022, configurando uma instabilidade estrutural que afeta toda a literatura baseada em dados dessa fonte.

\subsection{Limitações Relativas à Representatividade do Corpus}
\label{subsec:lim_representatividade}

\textbf{Restrição a dois veículos editoriais.} O corpus é composto exclusivamente por publicações institucionais dos veículos InfoMoney e InvestingBrasil, o que circunscreve o sentimento capturado
ao discurso editorial — mediado por critérios jornalísticos de seleção e enquadramento (\textit{framing}) — e não ao sentimento difuso e espontâneo dos participantes do mercado. O índice resultante deve ser interpretado como medida do \textbf{sentimento editorial financeiro}, e não do sentimento do investidor em sentido amplo, conforme concebido por \citeonline{baker2007} e \citeonline{bollen2011}.

\textbf{Ausência de anotação por múltiplos avaliadores.} O \textit{gold standard} construído para avaliação comparativa das abordagens foi anotado por um único avaliador, inviabilizando o cálculo de métricas de concordância interanotadores (e.g., Kappa de Cohen). A variabilidade intrínseca à tarefa de anotação de sentimento financeiro — especialmente para publicações de polaridade ambígua ou contextualmente dependente — representa uma fonte de incerteza não quantificada nos resultados de avaliação.


\subsection{Limitações Relativas aos Modelos Pré-treinados}
\label{subsec:lim_modelo}

\textbf{Tamanho reduzido do corpus de \textit{fine-tuning} do FinBERT-PT-BR.} Conforme reportado por \citeonline{santos2023finbert}, o \textit{fine-tuning} do FinBERT-PT-BR foi realizado sobre apenas 503 textos anotados manualmente, um volume consideravelmente menor do que o utilizado em modelos análogos para o inglês (e.g., o FinBERT original de \citeonline{araci2019}). Isso pode limitar a robustez do modelo em subdomínios ou estilos textuais subrepresentados no corpus de \textit{fine-tuning}.

\textbf{Possível desalinhamento de domínio do FinBERT-PT-BR.} O corpus de pré-treinamento do FinBERT-PT-BR é composto majoritariamente por notícias e relatórios financeiros, enquanto as publicações do X/Twitter apresentam características estilísticas distintas: linguagem mais informal, alta densidade de abreviações, emojis e \textit{hashtags}. Esse desalinhamento de domínio (\textit{domain shift}) pode degradar o desempenho do modelo em relação ao reportado nos experimentos originais. Essa limitação é parcialmente endereçada pela inclusão da abordagem BERTimbau ajustada (\autoref{subsec:bertimbau}), cujo \textit{fine-tuning} é realizado diretamente sobre o corpus de \textit{tweets} financeiros coletados, alinhando o modelo ao gênero textual e ao subdomínio de interesse.

\textbf{Cobertura léxica limitada no vocabulário financeiro especializado.} O OpLexicon v3.0 foi construído a partir de um corpus de propósito geral, com cobertura parcial do vocabulário técnico do mercado financeiro brasileiro — siglas de ativos, nomes de emissores, jargão de derivativos e expressões idiomáticas do setor não estão contemplados no léxico. Na prática, isso implica alta taxa de \textit{tokens} sem correspondência durante a pontuação agregada, reduzindo a magnitude dos escores e aumentando a probabilidade de classificação como Neutro. Essa limitação estrutural da abordagem léxica é um dos fatores que justificam a expectativa de que o FinBERT-PT-BR apresente desempenho superior nos resultados da avaliação comparativa.

\textbf{Limitações específicas da abordagem BERTimbau ajustada.} A estratégia de \textit{fine-tuning} local introduz um conjunto próprio de restrições metodológicas, distintas das limitações do FinBERT-PT-BR:

\begin{itemize}
	\item \textbf{Tamanho modesto do conjunto de \textit{fine-tuning} e variância das estimativas.} O corpus rotulado disponível para treinamento é composto por poucas centenas a poucos milhares de exemplos, volume que produz estimativas de desempenho com variância potencialmente elevada quando obtidas de uma única partição \textit{hold-out}. Embora o protocolo adotado --- particionamento estratificado 70/15/15 com semente fixa (\autoref{subsec:particionamento}) --- garanta reprodutibilidade, ele não quantifica a estabilidade das métricas em relação a diferentes realizações do particionamento. Uma análise complementar de robustez por validação cruzada estratificada de $k$ dobras ($k = 5$) é sugerida como extensão metodológica no \autoref{Cap:Ava_Exp};
	
	\item \textbf{Propagação do viés do anotador único.} A limitação documentada na \autoref{subsec:lim_representatividade} --- ausência de anotação por múltiplos avaliadores --- adquire gravidade adicional no contexto do BERTimbau ajustado: enquanto para o FinBERT-PT-BR e as abordagens léxicas o viés do anotador afeta apenas a métrica de avaliação, para o BERTimbau ele contamina também os pesos aprendidos pelo modelo. Decisões de polaridade subjetivas ou inconsistentes do anotador são internalizadas como sinal de treinamento, potencialmente amplificando padrões de erro sistemáticos na inferência;
	
	\item \textbf{Risco residual de \textit{overfitting}.} Apesar das três salvaguardas empilhadas --- \textit{early stopping} monitorando F1-Score macro na validação (\textit{patience}~=~2), decaimento de pesos (\textit{weight decay}~=~$10^{-2}$) e seleção automática do \textit{checkpoint} de melhor desempenho na validação ---, o risco de ajuste excessivo à distribuição de treinamento não é eliminado, especialmente se o corpus rotulado apresentar baixa diversidade estilística ou concentração temporal de eventos.
\end{itemize}

\subsection{Limitações de Escopo Analítico}
\label{subsec:lim_escopo}

\textbf{Ausência de modelagem preditiva.} Este trabalho não inclui análise de correlação ou modelagem preditiva entre o índice de sentimento construído e variáveis de mercado (retornos, volume
de negociação ou volatilidade). A validação do índice é exclusivamente qualitativa e descritiva. A investigação de relações quantitativas entre sentimento e comportamento de mercado, conforme documentado na literatura internacional \cite{tetlock2007,bollen2011,ranco2015}, constitui uma extensão natural deste trabalho, sugerida como direcionamento de pesquisa futura no \autoref{Cap:Conclusoes}.

\textbf{Ausência de análise de sentimento em nível de aspecto.} A abordagem adotada classifica cada publicação com uma única polaridade global (\textit{document-level}), sem distinguir o sentimento associado a entidades ou atributos específicos mencionados no texto (e.g., sentimento sobre determinada empresa, taxa de juros ou setor da economia). Conforme discutido na \autoref{subsec:abordagem_lexica}, a análise de sentimento baseada em aspectos (\textit{aspect-based sentiment analysis}) representa um nível de granularidade mais refinado, cuja implementação
demandaria recursos de anotação e infraestrutura computacional além do escopo deste trabalho.

O reconhecimento dessas limitações não invalida os resultados obtidos, mas circunscreve o alcance das generalizações admissíveis e orienta a agenda de pesquisa futura apresentada no \autoref{Cap:Conclusoes}.
